\documentclass[11pt,letterpaper]{article}
\usepackage[T1]{fontenc}
\usepackage[utf8]{inputenc}
\usepackage{lmodern}
\usepackage[margin=1in,headheight=15pt]{geometry}
\usepackage{amsmath,amssymb,amsthm,mathtools}
\usepackage{booktabs,longtable,array,tabularx}
\usepackage{ragged2e}
\usepackage{enumitem,xcolor,microtype,fancyhdr}
\usepackage{listings}
\usepackage{xurl}
\usepackage[unicode,breaklinks,colorlinks,linkcolor=blue!45!black,urlcolor=blue!50!black,citecolor=blue!45!black]{hyperref}
\hypersetup{pdftitle={Computational Complexity in Intuitionistic Primitive Recursive Arithmetic},pdfauthor={Research prepared for Vladimir Reshetnikov}}
\setlength{\parskip}{4pt}
\setlength{\parindent}{1.2em}
\setlist{itemsep=3pt,topsep=5pt}
\emergencystretch=2em
\pagestyle{fancy}
\fancyhf{}
\fancyhead[L]{\small Computational complexity in iPRA}
\fancyhead[R]{\small 20 November 2026}
\fancyfoot[C]{\thepage}
\renewcommand{\headrulewidth}{0.3pt}
\newtheorem{theorem}{Theorem}[section]
\newtheorem{lemma}[theorem]{Lemma}
\newtheorem{proposition}[theorem]{Proposition}
\newtheorem{corollary}[theorem]{Corollary}
\theoremstyle{definition}
\newtheorem{definition}[theorem]{Definition}
\newtheorem{remark}[theorem]{Remark}
\newcommand{\N}{\ensuremath{\mathbb N}}
\newcommand{\ipra}{\ensuremath{\mathsf{iPRA}}}
\newcommand{\PR}{\ensuremath{\mathsf{PR}}}
\newcommand{\Pclass}{\ensuremath{\mathsf P}}
\newcommand{\NP}{\ensuremath{\mathsf{NP}}}
\newcommand{\EXP}{\ensuremath{\mathsf{EXPTIME}}}
\newcommand{\PSPACE}{\ensuremath{\mathsf{PSPACE}}}
\newcommand{\NPSPACE}{\ensuremath{\mathsf{NPSPACE}}}
\newcommand{\NL}{\ensuremath{\mathsf{NL}}}
\newcommand{\Lclass}{\ensuremath{\mathsf L}}
\newcommand{\coNP}{\ensuremath{\mathsf{coNP}}}
\newcommand{\SAT}{\ensuremath{\mathsf{SAT}}}
\newcommand{\QBF}{\ensuremath{\mathsf{TQBF}}}
\newcommand{\Acc}{\operatorname{Acc}}
\newcommand{\Run}{\operatorname{Run}}
\newcommand{\Step}{\operatorname{Step}}
\newcommand{\len}[1]{\lvert #1\rvert}
\newcommand{\poly}{\operatorname{poly}}
\newcommand{\bit}{\operatorname{bit}}
\newcommand{\dotminus}{\mathbin{\dot{-}}}
\newcommand{\All}{\operatorname{All}}
\newcommand{\Any}{\operatorname{Any}}
\renewcommand{\Check}{\operatorname{Check}}
\newcommand{\code}[1]{\texttt{\small #1}}
\lstset{basicstyle=\small\ttfamily,columns=fullflexible,keepspaces=true,breaklines=true,frame=single,framerule=0.2pt,rulecolor=black!20,backgroundcolor=black!2,showstringspaces=false}
\begin{document}
\hypersetup{pageanchor=false}
\begin{titlepage}
\vspace*{0.45in}
{\small\sffamily FOUNDATIONS / CONSTRUCTIVE MATHEMATICS}\par
\vspace{0.35in}
{\Huge\bfseries Computational complexity\par in intuitionistic primitive\par recursive arithmetic\par}
\vspace{0.3in}
{\Large Bounded machines, explicit reductions, finite certificates,\par and a path to object-theory Lean proofs}\par
\vspace{0.45in}
{\large Research prepared for Vladimir Reshetnikov}\par
{\large 20 November 2026}\par
\vspace{0.45in}
\begin{abstract}
We develop a first-order, constructive foundation for complexity theory over
intuitionistic primitive recursive arithmetic. Binary words, machines, clocks,
finite computation histories, witnesses, circuits, and graphs are natural-number
codes. All inductions are reduced to quantifier-free induction by explicit
Boolean compilation of bounded invariants. Separate bit-machine cost proofs
prevent primitive-recursive definability from being mistaken for feasibility.

The development includes clocked presentations of the basic time and space
classes, polynomial reductions, witness-preserving Cook--Levin completeness,
constructive finite graph algorithms, Savitch simulation, nondeterministic
space complementation, quantified Boolean formula completeness, and a direct
diagonal language in exponential time with an explicit counterexample to every
polynomial descriptor. Further finite developments cover circuit counting,
randomized computation and amplification, counting witnesses, and the
polynomial hierarchy. The archive contains a kernel-checked Lean semantic
core, executable reference algorithms, and a detailed plan for lowering the
main results to the repository's actual \code{Prf PRAAx} judgment.

\textbf{Verification boundary.} The mathematical arguments are developed here;
the complete object-language \ipra{} derivations of Cook--Levin, Savitch,
space complementation, and the separation theorem have not been mechanically
constructed in this package. Checked Lean lemmas and remaining obligations
are identified separately.
\end{abstract}
\vfill
\noindent\small Repository basis: ProveIt commit
\texttt{83421593771f877bc1ae98f4229eb35507bff05b}.\par
\noindent\small This is a self-contained reconstruction and formalization foundation;
the classical complexity results are attributed, not claimed as new discoveries.
\end{titlepage}
\hypersetup{pageanchor=true}
\pagenumbering{roman}
\tableofcontents
\clearpage
\pagenumbering{arabic}
\section{Scope and the exact foundational claim}

Complexity theory studies the resources required by finite computations. Its
basic objects therefore fit particularly well into a finitistic arithmetic,
provided that two issues are kept separate: whether a function is available in
the theory, and how efficiently an ordinary bit machine computes it.
Exponentiation and exhaustive finite search are available in \ipra; neither is
declared to be a single physical machine step.

The linked overview \cite{WikiComplexity} supplies the subject map. We focus on
the standard discrete, worst-case theory: finite binary inputs, time, work
space, deterministic and nondeterministic machines, reductions, completeness,
circuits, and finite randomness. No assertion about arbitrary real-number
computation, quantum amplitudes, or every conceivable resource measure follows
from this scope. The distinction between intuitionistic bounded arithmetic and
\ipra{} also matters: Buss's theories characterize feasible functions under
specific definability conditions \cite{Buss1986}; \ipra{} itself contains
functions growing far faster than any polynomial.

\subsection{What the results mean}

There are three levels of statement in this report.
\begin{enumerate}
\item \emph{Arithmetic theorems}: first-order formulas about numerical codes,
proved by intuitionistic reasoning and permitted induction. These are the
actual target of the development.
\item \emph{Class notation}: a convenient external description of the
arithmetic theorems. For example, \(\NP\subseteq\PSPACE\) denotes a compiler
from each coded verifier with polynomial bounds to a coded polynomial-space
decider, together with an arithmetic correctness theorem.
\item \emph{Semantic adequacy}: in the ordinary standard natural numbers,
these presentations describe the familiar complexity classes. This is a
metamathematical identification, not an internal truth predicate for all
arithmetical formulas or all subsets of words.
\end{enumerate}

All classical machine and completeness facts used below are reconstructed with
explicit finite operations. Their \ipra{} status comes from those constructions
and their induction audits, not merely from the fact that they are standard
theorems. The supplied Lean project checks a useful part of the semantic
mathematics. A separate, explicit lowering step remains necessary to obtain
closed \code{Prf PRAAx} proofs for the main complexity theorems.

\begin{center}
\begin{tabularx}{\textwidth}{@{}p{0.31\textwidth}X@{}}
\toprule
Development & Constructive content supplied here\\
\midrule
Bounded arithmetic & Boolean quantifier elimination; finite least witnesses;
admissibility of bounded induction.\\
Machines and reductions & Total bounded simulation; explicit clocks; cost
accounting; reduction composition.\\
Time classes & \(\Pclass\subseteq\NP\subseteq\EXP\); closure constructions;
equivalence of verifier and branch presentations.\\
Completeness & Witness-preserving reductions to circuits and 3SAT;
\(\SAT\in\Pclass\leftrightarrow\Pclass=\NP\); \QBF{} completeness.\\
Space classes & \(\Lclass\subseteq\NL\subseteq\Pclass\subseteq\NP
\subseteq\PSPACE\subseteq\EXP\), \(\NPSPACE=\PSPACE\), and
\(\NL=\mathsf{coNL}\).\\
Separation & A fixed exponential-time Boolean language and a primitive-recursive
mismatch word for every polynomial machine descriptor.\\
Other foundations & Finite witness counting, fixed levels of the polynomial
hierarchy, rational probability and error amplification, finite circuit lower
bounds.\\
\bottomrule
\end{tabularx}
\end{center}

\subsection{Repository basis}

The repository defines \ipra{} as first-order intuitionistic logic with equality,
primitive-recursive function symbols and defining equations, successor axioms,
and quantifier-free induction \cite{ProveItPR}. Its maintained derivation
library has actual object-theory endpoints for bounded formulas, bounded least
witnesses, and bounded induction \cite{ProveItBounded,ProveItLeast,ProveItInd}.
These are directly relevant infrastructure. We inspected their source at the
commit on the title page. We did not rebuild the upstream project, and do not
inherit a research packet's verification claims without a separate proof.

\section{The arithmetic kernel and finite reasoning}
\label{sec:kernel}

\subsection{Language, axioms, and inference}

There is one numerical sort. Terms are built from variables, \(0\), successor
\(S\), and a symbol for each specified primitive-recursive definition. The
definitions are generated by zero, successor, projections, composition, and
primitive recursion. For
\[
 h(0,\vec x)=f(\vec x),\qquad
 h(Sn,\vec x)=g(n,h(n,\vec x),\vec x),
\]
both equations are axioms. Composition and projections have their analogous
defining equations. The successor axioms are
\[
 Sx\ne0,\qquad Sx=Sy\longrightarrow x=y.
\]
For each quantifier-free formula \(A(n,\vec x)\), the induction axiom is
\[
 A(0,\vec x)\longrightarrow
 \bigl(\forall n(A(n,\vec x)\longrightarrow A(Sn,\vec x))\bigr)
 \longrightarrow\forall n\,A(n,\vec x).
\]
The logic is intuitionistic first-order natural deduction: introduction and
elimination for \(\land,\lor,\to,\bot,\forall,\exists\), with the usual
eigenvariable conditions, equality reflexivity and substitution. Negation is
\(A\to\bot\). We use neither general excluded middle, double-negation
elimination, Markov's principle, choice, nor reflection of the theory's own
proofs. Equality of natural numbers and all explicitly decidable finite tests
do admit case distinction.

The symbol for a primitive-recursive definition is not a unit-cost operation
in the complexity model. Nor is there a total primitive-recursive evaluator
for all primitive-recursive syntax. If such a numerical evaluator \(U(e,x)\)
enumerated every unary primitive-recursive function, then
\(x\mapsto U(x,x)+1\) would be primitive recursive and diagonalize against its
own index. We only use fixed functions or a universal \emph{clocked Turing
machine} simulator, which has a different type of bound.

\subsection{Concrete numerical codes}

A binary word \(w=b_0\cdots b_{n-1}\) has code
\[
 \ulcorner w\urcorner=2^n+\sum_{i<n}b_i2^i.
\]
The marker \(2^n\) retains leading and trailing zero bits. The empty word has
code 1; code 0 is reserved as invalid and all language evaluators reject it.
The word length is the unique \(n\) with \(2^n\le x<2^{n+1}\), found by a
bounded search. We distinguish word length \(\len{x}\) from the numerical
magnitude of its code. Concatenation, extraction, reversal, and the formation
of unary blocks are primitive recursive.

For arrays with length \(r\) and entries less than \(B\), take \(B\ge2\) and
\[
 [a_0,\ldots,a_{r-1}]_B=\sum_{i<r}a_iB^i<B^r,
 \qquad a_i=\left\lfloor A/B^i\right\rfloor\bmod B.
\]
The length and base are supplied separately or included in a record header.
Tuples of bounded fields use mixed-radix coding. Unbounded finite records use
binary fields with self-delimiting lengths. An elementary numerical pairing,
for arithmetic purposes, is
\(\langle a,b\rangle=2^a(2b+1)-1\), whose projections use bounded division and
search. We do not assume this particular pairing is a length-efficient
serialization for reductions; those use explicit binary records. Every array
operation used in a proof has a numerical length/base bound. Huge bounds are
permissible for an arithmetic proof, but cannot be discarded in a machine-cost
argument.

\subsection{Boolean finite quantifiers and witness extraction}

Booleans are natural numbers 0 and 1. Their operations are primitive
recursive: \(\neg b=1\dotminus b\), \(a\wedge b=ab\), and
\(a\vee b=\min(1,a+b)\). For a fixed Boolean primitive-recursive predicate
\(b(i,\vec x)\), define
\begin{align*}
 \Any_b(0,\vec x)&=0,&
 \Any_b(r+1,\vec x)&=\Any_b(r,\vec x)\vee b(r,\vec x),\\
 \All_b(0,\vec x)&=1,&
 \All_b(r+1,\vec x)&=\All_b(r,\vec x)\wedge b(r,\vec x).
\end{align*}

\begin{lemma}[Bounded decisions and least witnesses]
\label{lem:finite}
For a fixed Boolean primitive-recursive \(b\), \ipra{} proves
\[
 \Any_b(r,\vec x)=1\ \leftrightarrow\ \exists i<r\ b(i,\vec x)=1,
 \qquad
 \All_b(r,\vec x)=1\ \leftrightarrow\ \forall i<r\ b(i,\vec x)=1.
\]
There is a primitive-recursive \(\mu_b(r,\vec x)\le r\), returning the least
successful index when one exists and the sentinel \(r\) otherwise. In
particular, bounded existence supplies an actual numerical witness.
\end{lemma}

\begin{proof}
Maintain a pair consisting of a found-bit and an index. Initially the bit is
0 and the index is 0. At stage \(i\), if no earlier success was found and
\(b(i,\vec x)=1\), set the found-bit to 1 and the index to \(i\); otherwise
leave the pair alone. After \(r\) stages, output the saved index when found,
and \(r\) otherwise. This is primitive recursion on a coded pair.

For clarity about the restricted induction, use the following elementary
prefix invariants: the found-bit equals \(\Any_b(i,\vec x)\); a found index
is below \(i\) and satisfies \(b=1\); and a non-found index means every
\(j<i\) has \(b(j,\vec x)=0\). The latter assertion is represented directly
by the finite Boolean product of \(1-b(j,\vec x)\). Preservation follows
from the two branches of the update and the defining fold equations. Each
resulting induction assertion is an equality of primitive-recursive terms,
or a quantifier-free combination of such equalities. For minimality, also
maintain that the finite product of failure bits below the saved index is 1.

One may recover the displayed quantified statements without induction on
them: a concrete successful index forces the OR fold to be 1, by the fold
equations and prefix splitting; if the fold is 1 the saved index gives the
required witness. If the AND fold is 1 each of its finitely many factors is
1. Conversely, a zero AND fold produces a least zero factor by the same
search construction. A premise that all factors are 1 excludes that case.
All implications are intuitionistically valid.
\end{proof}

\begin{theorem}[Bounded formula compilation]
\label{thm:compile}
For every fixed formula \(A(\vec x)\) whose quantifiers are bounded by
primitive-recursive terms, there is a Boolean primitive-recursive function
\(\chi_A(\vec x)\) and an \ipra{} derivation of
\[
 A(\vec x)\leftrightarrow\chi_A(\vec x)=1.
\]
Consequently induction for each such fixed bounded formula is admissible in
\ipra.
\end{theorem}

\begin{proof}
Equality and order have primitive-recursive characteristic bits. Conjunction,
disjunction and negation use the Boolean operations above. Since the already
compiled subformulas are decidable, implication is represented by
\(\neg\chi_A\vee\chi_B\); a four-case Boolean proof gives its equivalence
to \(A\to B\). Bounded existential and universal quantification use
Lemma~\ref{lem:finite}, with the other variables as parameters. This is a
structural induction in the metatheory over the fixed finite formula, giving
a specific function symbol and a specific derivation for that formula. It
does not assert a total internal truth evaluator for arbitrary formula codes
with arbitrary primitive-recursive symbols.

For bounded induction, translate the base case and successor implication to
the quantifier-free equation \(\chi_A(n,\vec x)=1\). Apply the permitted
induction axiom to this equation and translate back. The surrounding context
may contain other formulas; the formula being inducted on is still the
displayed equation.
\end{proof}

\subsection{How a finite proof is lowered}

An algorithmic proof below is admitted into \ipra{} by the following
discipline. Specify every constructor by composition, primitive recursion,
or explicitly bounded search. Give numerical bounds for the finite records
and quantifiers in its invariants. Prove base and successor steps by the
defining equations, finite Boolean case distinctions, and previously
established arithmetic lemmas. Apply Theorem~\ref{thm:compile} to the invariant.
Finally universally generalize the numerical parameters. Thus the universal
quantifiers in a final theorem do not become illicit unbounded quantifiers
inside the induction formula.

For example, ``all cells in the first \(t\) rows are correct'' is represented
by a bounded conjunction over row and cell indices. ``Every branch up to this
clock has the prescribed property'' is a bounded conjunction over branch
codes. When existence is required, a named bounded search or a named witness
constructor accompanies the statement. This discipline does not imply that
every true bounded formula is provable, or that every arbitrary Lean proof
has an \ipra{} translation. It specifies a concrete class of proof steps
and the translation of each step.

An especially useful consequence is that a finite search can be constructive
while taking exponential or much larger time. The distinction between
\emph{decidable in arithmetic} and \emph{decidable in polynomial time} is
preserved throughout the foundation.

\section{Machines, clocks, and resource accounting}
\label{sec:machines}

\subsection{A concrete bounded computation predicate}

For time we use ordinary deterministic single-tape machines, with one fixed
global alphabet \(\{0,1,\#,\mathrm{blank}\}\), a finite state set, an initial state,
and distinct accepting and rejecting states. The tape is indexed by integers;
the input occupies positions \(0,\ldots,n-1\), the remaining cells are blank,
and the head starts at 0. A transition reads a symbol, changes state, writes
a symbol, and moves by \(-1,0\), or 1. Halting states are made absorbing,
with no change to the tape or head. Undefined or malformed transitions reject.
A malformed machine code denotes a fixed rejecting machine. Numerical code
0 for a machine description is normalized to the empty binary description
(word code 1), which is interpreted as rejecting. Thus all coded
operations are total numerical functions.

A configuration consists of a state, an integer head position, and a finite
record of the nonblank tape cells. Encode a signed position by its sign and
absolute value, and normalize the tape record by removing blank boundary
cells. Reading an unrecorded cell returns blank. Lookup, replacement, adding
one new boundary cell, and normalization are bounded finite loops on this
record. Consequently \(\Step(e,C)\) and the initial configuration
\(\operatorname{Init}(e,x)\) are primitive recursive uniformly in the machine
description \(e\). Define
\[
 \Run(e,x,0)=\operatorname{Init}(e,x),\qquad
 \Run(e,x,t+1)=\Step(e,\Run(e,x,t)),
\]
and let \(\Acc(e,x,t)\) be 1 precisely when the resulting state is accepting.
For invalid input code 0 it is defined to be 0.

The equality \(\Acc(e,x,t)=1\) means ``accepts by time \(t\).'' It does not
assert that the original unclocked machine always halts. A clocked evaluator
runs the indicated computation and returns 0 on timeout, thereby always
producing a bit.

\begin{lemma}[Bounded computation]
\label{lem:boundedrun}
The functions \(\Run\) and \(\Acc\) are primitive recursive. Their defining
equations, trace composition, uniqueness of a checked deterministic trace,
and monotonicity of accepting-by-time are provable in \ipra.
\end{lemma}

\begin{proof}
Iteration is primitive recursion on \(t\), with \(e,x\) as parameters.
After \(t\) steps the head lies in \([-t,t]\), and all visited or initially
nonblank cells lie in \([-t,n+t]\). A window of
\(W=n+2t+3\) cells therefore contains the entire relevant tape, including
blank sentinels. For a fixed finite alphabet of size \(g\) and \(q\) states,
a window configuration is bounded, for example, by
\(qWg^W\), after choosing an offset for the head. Uniformly in machine codes,
the decoded numbers \(q,g\) provide the same primitive-recursive expression.

For composition prove
\(\operatorname{Iter}_{e}(r+s,C)=
\operatorname{Iter}_{e}(s,\operatorname{Iter}_{e}(r,C))\)
by induction on \(s\), using equality of coded configurations as the
induction formula. A trace checker tests its initial row and each adjacent
pair by a bounded conjunction. Inducting on the row number proves that every
checked row equals the corresponding iterate; the prefix assertion is a
finite conjunction and is Boolean-compiled. Finally an accepting
configuration is a fixed point of \(\Step\), so iteration preserves it.
This gives accepting-by-time monotonicity.
\end{proof}

For \(T\) steps, an entire window trace can be bounded by
\(B^{T+1}\), with \(B=1+qWg^W\), and can be generated by a single loop
accumulating rows. Such a trace is useful for correctness certificates; it is
not necessarily the storage used by the machine being analyzed.
When \(\Run\) uses the normalized finite-support representation, first
convert each row to the fixed window selected by \(T\). Bounded cell lookup
defines this conversion; the head/window bound proves that applying a step
then converting agrees with the fixed-window step throughout the trace.
The bound \(B^{T+1}\) applies to these converted rows, not to an unrelated
sparse-record serialization.

\subsection{Physical time is a second proof}

The arithmetic definition of \(\Run\) proves totality of the clocked
operation. It does not yet prove a feasible simulation. We use the following
coarse implementation bound; the later diagonal section gives its sweep
construction explicitly.

\begin{lemma}[Coarse universal implementation]
\label{lem:coarseU}
There is a fixed ordinary machine and a fixed positive integer \(C\) which
compute \(\Acc(e,x,T)\), with a unary clock supplied, in at most
\[
 C(\len e+\len x+T+1)^4
\]
steps. The transition-table simulation, its correctness, and the polynomial
overhead bounds have bounded invariants and \ipra{} proofs.
\end{lemma}

The constant \(C\) depends on the chosen universal implementation; this
package gives a constructive implementation recipe, not a numerically
optimized transition table for that universal machine. The full Lean
formalization must realize this recipe and establish its cost lemma. The
degree four is a deliberately generous bound for fixed-track sweep
simulation and table lookup, not a sharp time-hierarchy claim.

Fixed finitely many tapes can be represented on tracks of one tape, with
marked head positions. A complete sweep reads the marked symbols, and a
second updates the tracks and their markers. The resulting larger, but
fixed, tracking alphabet is encoded into the chosen global alphabet by
constant-length blocks, adding only a fixed multiplicative overhead.
At time \(t\), the occupied
interval has length at most a fixed multiple of \(n+t+1\). Thus \(T\)
multitape moves can be simulated in polynomial time. Conversely, one tape is
a special case of finitely many tapes. This supplies the model invariance
needed for polynomial and exponential classes. We invoke neither a
Church--Turing thesis nor a unit-cost model of arbitrary-precision arithmetic
as a resource theorem.

\subsection{Clock families and asymptotic bounds}

Write
\[
 p_{c,k}(n)=c(n+1)^k\quad(c\ge1,\ k\ge0),\qquad
 e_{c,k}(n)=2^{c(n+1)^k}.
\]
The evaluations are primitive recursive jointly in \(c,k,n\). Coefficients
and exponents are numerical data, with no arbitrary-function evaluation.
Whenever input reading is needed, enlarge a bound to dominate \(n+1\).
Every polynomial with nonnegative integer coefficients is bounded by some
\(p_{c,k}\). Sums and compositions have explicit majorants: for
\(p=a(n+1)^r\), \(q=b(n+1)^s\), with \(r,s\ge1\),
\[
 p(n)+q(p(n))+n+1
 \le\bigl(a+b(a+1)^s+1\bigr)(n+1)^{\max(r,rs,1)}.
\]
This follows by \(p(n)+1\le(a+1)(n+1)^r\), multiplication, and monotonicity
of powers. Any fixed polynomial simulation overhead can then be absorbed by
multiplying the degree and enlarging the coefficient.

It is useful to encode a polynomial as a finite list of nonnegative
coefficients as well. Its evaluation, addition, multiplication, and
substitution are primitive recursive, and their identities are proved by
finite sums and products. The Lean core implements an equivalent syntax of
constants, the variable, addition, and multiplication. The choice between
these polynomial representations is an implementation choice; it introduces
no new arithmetic strength.

Big-O notation can be expressed by explicit constants and a threshold.
For a known clocked algorithm, finitely many lengths below the threshold can
be covered by a bounded maximum and a larger coefficient. We normally
present bounds valid at every length, which avoids an unnecessary
eventual-domination obligation. For an arbitrary potentially partial machine,
one cannot form its worst-case halting time without a supplied halting bound.

\subsection{Separate conventions for space and output}

Space classes use a read-only input tape with endmarkers and a fixed number
of one-way work tapes. The input cells are not charged as work space; input
head positions still contribute to configuration counts. A cap \(s\ge1\)
allows cells \(0,\ldots,s-1\) on each work tape. An attempted transition
outside the cap has no accepting continuation. Section~\ref{sec:space}
constructs the finite configuration graph and proves cutoff and simulation
theorems. These explicit conventions are necessary for logarithmic space.

A function machine has a separate initially blank output tape, appending at
most one output bit per charged transition. Its output length is therefore
at most its clock, up to fixed record headers. A clocked transducer returns
the empty word on timeout; when used as a reduction, its entire clocked
behavior is the specified function. Output-size and runtime bounds are
stored together so that a later algorithm has a justified input-size bound.

\section{Class presentations and reductions}
\label{sec:classes}

\subsection{First-order definitions of P, NP, and their relatives}

In every language-equivalence and reduction statement, input variables range
over valid word codes \(x\ge1\). We write this guard explicitly in the
defining sentences below and leave it implicit in subsequent displays.
Transducers output valid word codes, including code 1 for the empty word.
The reserved code 0 is outside the language universe; in particular,
complementation is relative to valid words.

A deterministic polynomial descriptor is \(d=(e,c,k)\), with \(c\ge1\).
Malformed descriptors are normalized to a fixed rejecting descriptor;
the raw description component \(e=0\) is normalized as specified above.
Define the Boolean function
\[
 P_d(x)=\Acc(e,x,p_{c,k}(\len x)).
\]
A verifier descriptor \(v=(e,c,k,a,l)\) has witness length
\(p(n)=c(n+1)^k\) and time clock \(q(n)=a(n+1)^l\). Its input record is
\(x\#y\), where \(y\) has exactly \(p(n)\) bits and the separator is part
of the fixed finite machine alphabet. To keep binary word codes and this
initialized tape distinct, define \(\Acc_{\#}(e,x,y,t)\) by the same step
iteration as \(\Acc\), but starting with the tape record \(x\#y\).
This is a fixed primitive-recursive function of two binary word codes and
the clock. Its simulation cost uses tape-input length \(\len x+\len y+1\).
Equivalently, an ordinary binary machine can read a self-delimiting binary
pair and construct this record with polynomial overhead.
Let \(\operatorname{bits}_p(j)\)
denote the \(p\)-bit representation of \(j<2^p\), including leading zeros.
Set
\[
 N_v(x)=\bigvee_{j<2^{p(n)}}
 \Acc_{\#}(e,x,\operatorname{bits}_{p(n)}(j),q(n)),\qquad n=\len x.
\]
All operations in this definition are fixed primitive-recursive functions.
An exact-length certificate entails no loss: a certificate of length at most
\(p\) can be represented by a unary length delimiter, its contents, and
padding in a fixed length \(2p+1\) record.

\begin{definition}[Representing a language]
For a fixed arithmetic formula \(A(x)\), write
\[
 \operatorname{PRep}(d,A):=\forall x\ge1\bigl(A(x)\leftrightarrow P_d(x)=1\bigr),
\]
and define \(\operatorname{NPRep}(v,A)\) by replacing \(P_d\) with \(N_v\).
The abbreviations \(A\in\Pclass\), \(A\in\NP\) mean existence of the
corresponding numerical descriptor. This use of \(A\) is schematic; there is
no internal quantification over arbitrary predicates. Alternatively, work
entirely with the Boolean functions \(P_d,N_v\), and compare them by
\(\forall x\ P_d(x)=N_v(x)\).
\end{definition}

Exponential time uses \(\Acc(e,x,e_{c,k}(\len x))\). Space presentations
use acceptance in the capped configuration graph for an explicit
\(s(n)=c\lceil\log_2(n+2)\rceil\) or \(s(n)=p_{c,k}(n)\), as developed
below. Integer logarithms mean bounded search for a power of two, not real
analysis. Deterministic versus nondeterministic edges distinguish
\(\Lclass,\NL\) and \(\PSPACE,\NPSPACE\).

For coNP use \(1-N_v(x)\). Its direct logical form is a bounded universal
statement that every certificate is rejected. Each such input statement is
decidable in \ipra. This does not place its language in NP, because the
exhaustive decision algorithm can take exponential time.

\begin{theorem}[Verifier and branch presentations]
Polynomial-time nondeterministic acceptance and the verifier presentation
\(N_v\) agree, with primitive-recursive translations of descriptors and
bounded witness/branch correctness proofs in \ipra.
\end{theorem}

\begin{proof}
Normalize a fixed finitely branching transition table so that each choice
is coded by a fixed number \(b\) of bits. A branch of length at most \(T\)
is coded by \(bT\) choice bits, with padding after halting. A deterministic
verifier replays those choices, rejecting an invalid transition code. Its
simulation has polynomial overhead by the preceding sweep implementation.
The witness constructor records the choices of a supplied branch, and the
inverse replays them. The prefix invariant asserts equality of the recorded
and replayed configurations at every earlier time, a bounded formula.

Conversely, a nondeterministic machine guesses the prescribed polynomial
number of certificate bits, then runs the fixed clocked verifier. Its clock
includes guessing, input assembly and verification. The constructed branch
exists exactly when the bounded witness disjunction is 1. All loops and
record lengths are explicitly bounded. No infinite branch principle is used.
\end{proof}

\subsection{Elementary closure and exhaustive simulation}

\begin{theorem}[Closure and basic containments]
\label{thm:closure}
\ipra{} proves the numerical compiler theorems corresponding to:
\(\Pclass\subseteq\NP\); closure of P under complement, finite union and
intersection; closure of NP under finite union and intersection; and
\(\NP\subseteq\EXP\).
\end{theorem}

\begin{proof}
For \(\Pclass\subseteq\NP\), ignore the witness and run the supplied
deterministic descriptor. To complement P, evaluate the entire clocked
Boolean computation and negate its answer, including the timeout answer.
Merely interchanging the original machine's halting labels is insufficient
when the original machine can time out. Sequentially evaluate two clocked
descriptors and combine their bits by AND or OR to obtain intersection or
union. The resulting sums and simulation overhead have polynomial majorants.

For intersection of two NP presentations, concatenate their two
certificates and verify both. The forward witness map concatenates; the
reverse projects. For union use a selector bit and a certificate field
padded to the larger of the two prescribed lengths; the selector determines
which verifier is run. These are explicit witness maps, so the existential
equivalences are intuitionistic.
If a concatenated, padded, or selector-augmented certificate length is not
already a monomial, choose the explicit monomial majorant from
Section~\ref{sec:machines} and ignore the additional padding. This keeps the
constructor inside the displayed descriptor family.

Finally enumerate \(2^{p(n)}\) certificates, replaying the verifier for
\(q(n)\) steps on each. A fixed implementation takes at most
\[
 A\,2^{p(n)}(n+p(n)+q(n)+1)^b
\]
steps for fixed integers \(A,b\). Replace the polynomial factor and \(A\)
by powers of two using \(r\le2^r\), and absorb the resulting polynomial
exponent into \(2^{r_0(n)}\). The scan's OR accumulator is exactly the
defining bounded disjunction. Its induction invariant compares the
accumulator with the finite prefix OR. Exhaustive search here proves an
exponential upper bound, not a polynomial one.
\end{proof}

\subsection{Reductions carry algorithms and bounds}

\begin{definition}[Polynomial many-one reduction]
A reduction from Boolean presentations \(A\) to \(B\) is a clocked
polynomial transducer \(f\), together with
\[
 \forall x\quad A(x)=B(f(x)).
\]
For formula presentations replace equality by equivalence. Its certificate
records a polynomial time majorant and an output-length majorant. A
\emph{witness-preserving} reduction additionally specifies primitive-recursive
maps transporting source witnesses to target witnesses and target witnesses
back to source witnesses, with their correctness and size bounds.
\end{definition}

\begin{theorem}[Composition and downward closure]
\label{thm:reduction}
Polynomial many-one reduction is reflexive and transitive. If \(A\le_m^pB\)
and \(B\in\Pclass\), then \(A\in\Pclass\). If \(B\in\NP\), then
\(A\in\NP\). These statements have explicit descriptor compilers and
\ipra{} proofs.
\end{theorem}

\begin{proof}
The identity transducer gives reflexivity. For composition, run \(f\), then
\(g\) on its output. If the respective bounds are \(p(n)\) and \(q(m)\),
then \(\len{f(x)}\le p(n)\), so the useful work is at most
\(p(n)+q(p(n))\), with an explicit copying and simulation overhead.
The polynomial-majorant calculation in Section~\ref{sec:machines} supplies
a resulting clock. Correctness is ordinary equality transitivity:
\(A(x)=B(f(x))=C(g(f(x)))\).

For downward closure of P, compose with the target decider. For NP, compute
\(f(x)\), take a certificate for the target verifier at that input, and
verify it. If target witness length is \(r(m)\), use a fixed padded field
of length \(r(p(n))\), discarding padding after the required length
\(r(\len{f(x)})\). Its existence is equivalent by extension and projection.
The combined verification time is polynomial. No unbounded choice of a
machine or a witness occurs in these compiler proofs.
\end{proof}

A language is hard for a presentation class if every descriptor in the class
has such a reduction to it; completeness adds membership. In arithmetic this
can be written as a numerical quantification over descriptors, reduction
codes, and their pointwise correctness. The standard completeness proofs
below do more: they give a primitive-recursive constructor for the reduction
code rather than only asserting its existence.

\subsection{Search, functions, promises, and oracle queries}

The class FP is presented by polynomially clocked transducers. A search
problem has a Boolean polynomial-time relation \(R(x,y)\) and a polynomial
balance bound \(\len y\le p(\len x)\). A solution procedure returns a witness
when one exists, and a distinguished failure record otherwise. Total search
adds a proved assertion that a witness always exists; this assertion must
not be inferred just from the verifier's totality.

A promise problem is specified by a promise formula \(H(x)\), a yes-condition
and a no-condition, with the claimed algorithmic behavior required only under
\(H(x)\). Completeness arguments must preserve the promise, not just map
yes-instances correctly. A numerical certificate for a promise can be handled
as ordinary input data; an arbitrary truth oracle is not silently inserted.

A polynomial oracle reduction consists of a clocked query machine with
polynomially many polynomial-length queries. For a fixed Boolean oracle
presentation, its query answers can be substituted to obtain a primitive-
recursive bounded computation. If no primitive-recursive oracle definition
is supplied, use a finite transcript of queries and supplied answers and a
relational correctness statement. This gives an appropriate base for later
oracle complexity while respecting first-order arithmetic. No unrestricted
oracle truth predicate is part of the present theory.

\section{Witness-preserving NP completeness}
\label{sec:cook}

The classical starting point is Cook's theorem on polynomially bounded
nondeterministic computations and propositional problems \cite{Cook1971}.
Cook's original paper uses polynomial oracle reducibility and tautologies.
Here we prove a modern many-one, witness-preserving circuit and CNF version,
and make its arithmetic compilation explicit.

\subsection{Circuits and normalized SAT instances}

A Boolean circuit is a finite list of gates in topological order. Input
variables precede internal gates; every wire of a gate refers to an input or
an earlier gate. The allowed internal operations are constants, NOT, AND,
and OR; OR may instead be expanded using NOT and AND. A designated wire is
the output. The circuit code includes its number of inputs and its gate
list. The validity test checks all index bounds. Evaluation fills a Boolean
array from left to right, so it is primitive recursive. Ordinary sequential
bit operations evaluate it in polynomial time in its serialization length;
a simple implementation can scan the previous wire table for every lookup.

A CNF is a finite list of clauses, and a clause a finite list of signed
variable indices. For \(3\SAT\), clauses have at most three literals.
We use normalized instances: the number \(m\) of variables is at most the
encoding length, all indices lie in \(0,\ldots,m-1\), and the input is a
well-formed binary record. A malformed instance is unsatisfiable by
definition. General binary variable names can be renamed to the distinct
names appearing in the formula, in polynomial time by pairwise comparison
and scanning. This matters because the largest numerical identifier could
be exponentially larger than the formula length. The witness is one bit per
occurring variable, not one bit per number below that largest identifier.

Write \(\Check(F,a)=1\) for a well-formed assignment satisfying all clauses.
This is a bounded conjunction of bounded disjunctions of literal bits. It has
a polynomial bit-machine implementation. Thus normalized CNF-SAT and 3SAT
have NP presentations, with unused certificate bits ignored after the
assignment's required length.

\subsection{Unrolling a verifier locally}

Fix a verifier descriptor \(v\), input \(x\) of length \(n\), witness length
\(p=p(n)\), and clock \(T=q(n)\). Convert its fixed machine to the
single-tape model if necessary, enlarging its polynomial clock. Put
\(x\#y\) on the initial tape and use a window of
\[
 W=n+p+2T+4
\]
cells. The window contains the whole input and at least \(T+1\) blank cells
of margin on either side, with harmless surplus cells. No head can reach
the external boundary during the relevant updates.

A cell label is either an ordinary tape symbol or a pair consisting of a
state and a symbol, meaning that the head is at that cell. Thus the finite
label set is
\[
 \Delta=\Gamma\ \sqcup\ (Q\times\Gamma).
\]
For valid configurations, the label of cell \(j\) after one move depends
only on the previous labels at \(j-1,j,j+1\). To see this, if the head is
at \(j\) it determines the new written symbol and whether it remains; if
it is at a neighbor it may move into \(j\); if neither occurs the symbol is
unchanged and the cell has no head. There is a fixed finite function
\(F_v:\Delta^3\to\Delta\) implementing these cases. Its values on
impossible triples, such as several conflicting heads, can be chosen
arbitrarily. The induction below only reaches valid rows.

Encode labels by a fixed number of bits, or by one-hot bits. Each output
bit of \(F_v\) is a fixed finite Boolean function. Its truth-table DNF is
an explicit circuit: OR together the ANDs specifying exactly the input
tuples on which that bit is 1. The number of gates is a computable constant
depending on the finite machine \(v\).

Construct \(C_{v,x}\) in row order. At row 0 hardwire the input, separator,
blank cells and initial head; use the \(p\) certificate bits as input wires.
For every time \(t<T\) and every cell, insert a fresh copy of the local
update circuit wired to the previous three cell labels. In the final row,
OR the tests for an accepting head. Absorbing terminal states mean the final
row reports acceptance by \(T\), including earlier halting.

\begin{lemma}[Uniform circuit construction]
\label{lem:unroll}
There is a primitive-recursive circuit constructor satisfying
\[
 \operatorname{Eval}(C_{v,x},y)
   =\Acc_{\#}(e,x,y,T).
\]
For each fixed verifier descriptor its gate count is bounded by
\[
 G\le a_v(T+1)W,
\]
where \(a_v\) is an explicitly computable constant. Its serialization and
its generation take polynomial time in \(n\).
\end{lemma}

\begin{proof}
The constructor is a collection of bounded loops over rows, cells, label
bits and the fixed local truth table. Its intermediate gates are finite
arrays, and every predecessor index is calculated from earlier row and gate
indices. This supplies a primitive-recursive syntax constructor, including
its validity proof.

Induct on \(t\). The invariant says that, for all \(j<W\), the Boolean
label in row \(t\) equals the corresponding machine cell label. It also
records that the head is unique and stays inside the allocated window.
The initial row is the definition of the machine input. At the successor
step the three previous labels are correct, and the fixed truth table is
precisely the local transition. Each claim is a bounded finite conjunction;
Theorem~\ref{thm:compile} replaces it by a Boolean equation. The final
acceptance OR proves the displayed identity.

There are \(TW\) local blocks, plus \(O(W)\) initialization and output
gates, giving the count. Each wire index needs
\(O(\log(G+p+2))\) bits. A generator which repeatedly scans its counters
and previous output takes, for example, a fixed multiple of
\((G+p+n+1)^3\) bit operations. This intentionally loose bound follows
by counting the bounded output and index-manipulation loops. Since \(p,T\)
are fixed polynomials in \(n\), both output length and runtime are
polynomial. The polynomial constants may depend on the verifier, as
many-one hardness allows. A primitive-recursive algorithm for building the
local truth tables is uniform in the verifier code, without asserting that
this preprocessing is polynomial in that variable code.
\end{proof}

\subsection{Tseitin constraints and both witness directions}

For each internal gate introduce a fresh variable \(z\). Retain the input
variables. Use the following clauses:
\begin{center}
\begin{tabular}{@{}ll@{}}
\toprule
Gate equation & Clauses\\
\midrule
\(z=\neg a\) & \((z\lor a),\ (\neg z\lor\neg a)\)\\
\(z=a\land b\) & \((\neg z\lor a),\ (\neg z\lor b),\
 (z\lor\neg a\lor\neg b)\)\\
\(z=a\lor b\) & \((z\lor\neg a),\ (z\lor\neg b),\
 (\neg z\lor a\lor b)\)\\
\(z=1\), \(z=0\) & \((z)\), respectively \((\neg z)\)\\
\bottomrule
\end{tabular}
\end{center}
Append a unit clause asserting the output variable. Also append the
tautology \((y_i\lor\neg y_i)\) for each of the \(p\) original input
variables. These tautologies retain even unused witness bits under
occurrence-based variable normalization, without constraining their values.
There are at most \(3G+p+1\) clauses. All have width at most three. If exactly three literals
are required, repeat a literal in shorter clauses. Output length is
\(O((G+p+1)\log(G+p+2))\), and the gate-to-clause conversion is polynomial
time. A direct Boolean case analysis, at most eight cases, proves that each
gadget is equivalent to its gate equation, even if incoming wires coincide.

Let \(F_v(x)\) be the resulting CNF. Define
\(\operatorname{Ext}_v(x,y)\) to be the input bits followed by all evaluated
gate bits. Define \(\operatorname{Restr}_v(x,a)\) to be the assignment's
input-bit prefix. Both are primitive-recursive functions with polynomial
length and polynomial-time implementations for fixed \(v\).

\begin{theorem}[Witness-preserving Cook--Levin]
\label{thm:cook}
For each verifier descriptor \(v\), \ipra{} proves
\begin{align*}
 \Acc_{\#}(e,x,y,T)=1
 &\ \longrightarrow\
 \Check(F_v(x),\operatorname{Ext}_v(x,y))=1,\\
 \Check(F_v(x),a)=1
 &\ \longrightarrow\
 \Acc_{\#}(e,x,\operatorname{Restr}_v(x,a),T)=1.
\end{align*}
The construction is a polynomial many-one reduction from \(N_v\) to 3SAT.
The code constructor, clocks, witness maps, and the correctness derivations
are effective finite constructions, using only the arithmetic kernel above.
\end{theorem}

\begin{proof}
Forward, evaluated gate bits satisfy every defining gate equation and hence
every gadget. Lemma~\ref{lem:unroll} gives the output bit 1, satisfying the
final unit clause.

Backward, fix a satisfying assignment. In topological order, each gate's
clauses force its assigned value to equal the Boolean operation on its
assigned predecessors. By induction, these predecessor values already equal
the actual evaluation from the restricted input. Therefore the gate value
does too. Use as invariant the bounded conjunction over all gates already
processed. The output clause now forces the actual output to be 1, and
Lemma~\ref{lem:unroll} translates it to verifier acceptance. These arguments
do not infer a witness from a double negation; they construct and project
witnesses explicitly. The size and runtime bounds follow from the two
constructors and reduction composition.
\end{proof}

\subsection{The meaning of P versus NP inside the theory}

Choose a fixed normalized 3SAT verifier and let \(\operatorname{sat}(x)\)
be its exhaustive Boolean satisfiability function. The following are
legitimate first-order sentences:
\begin{align*}
 \mathsf{SatP}&:\ \exists d\ \forall x\ge1\ P_d(x)=\operatorname{sat}(x),\\
 \mathsf{PNP}&:\ \forall v\ \exists d\ \forall x\ge1\ P_d(x)=N_v(x).
\end{align*}
Theorem~\ref{thm:cook} and the reduction compiler prove
\(\mathsf{SatP}\leftrightarrow\mathsf{PNP}\). One direction instantiates
\(v\) with SAT. For the other, given a SAT descriptor \(d\), construct
for each \(v\) the decider obtained by compiling \(F_v\) and then running
\(d\). In fact this construction gives a primitive-recursive function
\(\operatorname{Solve}(d,v)\), so the implication includes explicit
uniform machine selection. Nothing here supplies the premise
\(\mathsf{SatP}\). P versus NP remains unresolved \cite{ClayPNP}.

There are several distinct constructive lower-bound formulations. The
negation of \(\mathsf{SatP}\) is
\(\forall d\,\neg\forall x(P_d(x)=\operatorname{sat}(x))\). The stronger
statement
\[
 \forall d\ \exists x\ P_d(x)\ne\operatorname{sat}(x)
\]
actually supplies a mismatch for each descriptor. Even though equality at
each fixed \(x\) is decidable, intuitionistic logic does not generally
convert a negated unbounded universal into an existential counterexample.
The diagonal theorem below deliberately proves the stronger form for its
exponential-time language, with a named primitive-recursive witness.

\subsection{Search and decision, conditionally}

Assume a fixed polynomial SAT decider. For a satisfiable formula with
\(m\) variables, query whether its restriction with variable 0 set to 0 is
satisfiable. Keep 0 if yes, otherwise keep 1. Continue through all variables.
At each stage the formula with the chosen prefix remains satisfiable:
the disjoint cases for the next Boolean variable exhaust its assignments,
and a negative answer to one leaves a witness in the other. After \(m\)
stages the accumulated assignment itself satisfies the original formula.
An initial satisfiability test returns failure for unsatisfiable formulas.
Use ordinary CNF-SAT here; it is also reducible to 3SAT by the already
proved verifier construction, with polynomial overhead.

The stage invariant is bounded satisfiability of a finite restriction, so
it compiles to quantifier-free induction. The procedure makes at most
\(m+1\) polynomial-length SAT queries and polynomial bookkeeping steps.
For an arbitrary polynomially balanced NP search relation, first use
Theorem~\ref{thm:cook}, find a satisfying assignment, and project it back.
Thus a SAT decider yields polynomial search algorithms with explicit
witness recovery. A converse follows by applying a total search-with-failure
solver to SAT and checking its output. These are conditional compiler
theorems, not an unproved equivalence between an arbitrary relation and an
arbitrarily supplied total function.

Complementing the pointwise reduction also proves UNSAT coNP-complete in
the corresponding presentation. To reduce formula tautology to CNF-UNSAT,
negate the formula and apply the Tseitin encoding to its evaluation circuit.
The resulting CNF is unsatisfiable exactly when the original formula is a
tautology. The reverse reduction to general formula tautology simply negates
the CNF. These finite Boolean operations
do not establish \(\NP=\coNP\).

\section{Space, reachability, and finite counting}
\label{sec:space}

Space complexity makes the distinction between an arithmetical definition and
an efficient implementation particularly clear. A large finite table gives a
convenient primitive-recursive definition of reachability. A different,
depth-first implementation establishes a small space bound. The arithmetic
proof must connect the two. We give these connections explicitly, and then
reconstruct nondeterministic complementation and quantified Boolean
completeness in the same finite setting.

\subsection{Configurations and their bounds}

Fix a machine with $q$ states, $k$ work tapes, and a work alphabet of cardinality
$g\geq 2$. The binary input tape is read-only and has endmarkers. For an input
of length $n$, restrict each work tape to cells $0,\ldots,s-1$, with $s\geq 1$.
A configuration consists of the state, the input-head position, all work-head
positions, and all work-tape contents. Consequently
\[
 N_0=q(n+2)s^k g^{ks}
\]
is an exact bound on syntactic configurations. Mixed-radix encoding gives
primitive-recursive encoding and projection functions. In particular, tape
contents alone would give an incorrect configuration count: head positions
and the state matter.

The configuration graph contains an edge precisely when a legal machine
transition connects the two configurations without exceeding the cap. Delete
moves outside the cap. Add one new target vertex, numbered $N_0$, replace
the outgoing edges of every accepting configuration by its sole edge to
that target, and give the target no outgoing edges. This preserves
acceptance and keeps deterministic configuration graphs deterministic.
Put $N=N_0+1$. Initial and
target vertices are explicit functions of the machine and input. In a padded
binary encoding, unused codes simply have no useful outgoing edges.

Let $\ell$ be the least integer with $2^\ell\geq N+1$. For a fixed machine,
\[
 \ell\leq C_M\bigl(s+\lceil\log_2(n+2)\rceil+1\bigr).
\]
Here and below logarithm notation abbreviates integer binary length, defined
by bounded search among powers of two. The bound follows from the product
formula and $\log_2 s\leq s$ for $s\geq1$; it involves no analysis of real
logarithms.

Acceptance in the capped graph means existence of an accepting branch that
stays inside the cap. This agrees with ordinary acceptance when all branches
of the original machine obey the stated space bound. Branches may loop. If a
physical simulation must compute a function $s(n)$, that computation must
also obey the claimed space bound. The logarithmic and polynomial caps used
below have explicit suitable constructors. We make no unrestricted claim
that computing an arbitrary primitive-recursive cap is inexpensive.

\subsection{Paths, cycle removal, and the deterministic cutoff}

A path of length $m$ is a coded sequence $v_0,\ldots,v_m$ of vertices below
$N$, with the required adjacent edges. Its code is bounded by $N^{m+1}$,
with the length supplied separately. The predicate
$\mathrm{PathLe}(a,b,m)$ is a bounded disjunction over lengths $r\leq m$ and
codes below $N^{r+1}$. This deliberately expensive definition immediately
establishes primitive recursiveness.

\begin{lemma}[Constructive cycle removal]
Given a path from $a$ to $b$ in an $N$-vertex graph, a primitive-recursive
operation produces a path with the same endpoints and length at most $N-1$.
\end{lemma}
\begin{proof}
If positions $i<j$ contain the same vertex, delete positions
$i+1,\ldots,j$. The edge after position $i$ remains legal because the old
vertices at $i$ and $j$ coincide. Endpoints are unchanged, and the length
drops by at least one. Find the least repeated pair by bounded search, and
perform this deletion at most $m$ times, doing nothing once there is no
repeat. This is a fixed bounded iteration. Its invariant says that endpoints
and adjacency are preserved and that either no repeat remains or the length
has decreased. After $m$ rounds no repeat can remain. A duplicate-free list
of vertices below $N$ has at most $N$ entries, by finite counting.
\end{proof}

All assertions in this proof are bounded properties of the original path and
its successive codes. Their characteristic functions turn the invariant into
a quantifier-free formula. The finite counting assertion can itself be proved
using indicator sums over the $N$ possible vertex labels.

For a deterministic capped machine, an accepting trace can therefore be
shortened to one with fewer than $N$ edges. Determinism ensures that this
short trace agrees with the initial segment of the unique run. Thus the
following bounded implication holds for every supplied $t$: acceptance by
$t$ implies acceptance before $N$. A capped recognizer can be made a decider
for its acceptance language by rejecting at this cutoff. There is no need to
prove an unbounded statement about every future configuration as an
induction invariant.

\subsection{A polynomial graph algorithm}

For source $a$, let $S_0(v)=[v=a]$ and define
\[
 S_{i+1}(v)=S_i(v)\mathbin\lor
       \bigvee_{u<N}\bigl(S_i(u)\mathbin\land E(u,v)\bigr).
\]
Encode $S_i$ as an $N$-bit integer. Updating the vector is a fixed
primitive-recursive operation, so recursion on $i$ defines every vector.
An induction shows that $S_i(v)=1$ precisely when there is a path of length
at most $i$ from $a$ to $v$: a newly marked vertex is reached by extending
an old path, and the penultimate vertex of any nonempty path supplies the
converse. Witness paths can be stored in the mathematical construction,
using bounded least search, even when the implementation needs only bits.

The conjunction of this invariant over $v<N$ is bounded and can be
Boolean-compiled before induction. Cycle removal shows that $S_{N-1}$ is
the full reachable set. Computing $N$ rounds by scanning all ordered pairs
uses $O(N^3)$ adjacency queries and polynomial time on an ordinary explicit
graph representation. For a logarithmic-space machine the configuration
count is polynomial in input length, and adjacency is polynomial-time
decidable. This gives an explicit proof of $\mathrm{NL}\subseteq\mathrm{P}$.

\subsection{Savitch's theorem with a primitive-recursive audit}

The classical square-space simulation is due to Savitch
\cite{Savitch1970}. Define Boolean tables
\begin{align*}
 H_0(a,b)&=[a=b]\lor E(a,b),\\
 H_{j+1}(a,b)&=\bigvee_{z<N}
               \bigl(H_j(a,z)\land H_j(z,b)\bigr).
\end{align*}
For a direct primitive-recursive definition, compute the entire $N$ by $N$
table and update it by this formula. Table codes are below $2^{N^2}$.
Recursion on a single table code is ordinary primitive recursion; an
unanalysed recursive subroutine is not being treated as a totality axiom.

\begin{theorem}[Finite repeated squaring]
For every $j$, $H_j(a,b)=1$ if and only if there is a path from $a$ to $b$
of length at most $2^j$.
\end{theorem}
\begin{proof}
The base case permits lengths zero and one. Concatenating two paths of
length at most $2^j$ gives a path of length at most $2^{j+1}$. Conversely,
in a path of length $m\leq2^{j+1}$, take the vertex at position
$\min(m,2^j)$. The two resulting segments both have length at most $2^j$.
This midpoint witnesses the displayed disjunction. Split and join operations
on the coded paths are primitive recursive.
\end{proof}

For the iPRA induction, bound path codes by $N^{2^j+1}$ and conjoin over
$a,b<N$. The resulting assertion is bounded, even though its bound is large.
Its characteristic function supplies the quantifier-free invariant.
If $j$ is the least integer with $2^j\geq N$, repeated squaring decides
reachability.

The small-space implementation evaluates this recurrence depth first. Each
frame stores $a,b$, the midpoint loop index, the current level, a phase flag,
and one child answer. Evaluate one child, retain only its Boolean result,
then evaluate the other. At most $j+1$ frames are active. Each frame uses
$O(\ell+\log(j+2))$ bits, and $j\leq\ell$, giving $O(\ell^2)$ bits in all,
plus reusable adjacency scratch space. The recursion tree has at most
\[
 1+(2N)+\cdots+(2N)^j\leq(2N+1)^{j+1}
\]
nodes. Multiplication by the explicit per-node bookkeeping and adjacency
cost gives a primitive-recursive clock for an iterative stack machine.

Quantifier-free induction on the level, bounded over endpoints and the
relevant stack codes, proves that this clocked evaluator returns the table
value and obeys its frame bound. This is the essential bridge between the
large table definition and the small-space implementation. The table alone
would not prove a space bound.

\begin{theorem}[Square-space simulation for explicit caps]
For the explicit machine and cap data above, nondeterministic acceptance is
decidable in deterministic space
$O((s+\log(n+2))^2)$. If $s\geq\lceil\log_2(n+2)\rceil$, the bound is
$O(s^2)$. In particular, $\mathrm{NPSPACE}=\mathrm{PSPACE}$.
\end{theorem}
\begin{proof}
Use the configuration graph and the depth-first evaluator. The vertex code
length is $O(s+\log(n+2))$. Testing one transition needs only comparable
scratch space. For a polynomial cap its square is polynomial; the reverse
containment uses a deterministic machine as a nondeterministic machine with
one possible transition.
\end{proof}

\subsection{Nondeterministic complementation by exact counting}

The counting method behind nondeterministic complementation is due
independently to Immerman and Szelepcs\'enyi; we cite Immerman's account
\cite{Immerman1988}. Here we give the finite certificate construction and
its restricted-induction audit.

Fix source $a<N$ and put
\[
 R_i(v)=[\mathrm{PathLe}(a,v,i)],\qquad
 c_i=\sum_{v<N}R_i(v).
\]
In particular $c_0=1$. We will compute $c_{N-1}$ nondeterministically so
that every surviving branch returns the correct count and at least one
branch survives.

\begin{lemma}[Exact enumeration]
Suppose the correct value $c_i$ is supplied. A nondeterministic scan using
$O(\log N)$ work bits, plus adjacency scratch space, can enumerate precisely
the vertices reachable within $i$ steps. Every successful scan is exact,
and some scan succeeds.
\end{lemma}
\begin{proof}
Visit vertices $u=0,\ldots,N-1$ in order and guess a mark $b_u$. If $b_u=1$,
guess a path from $a$ to $u$ of length at most $i$, checking one edge at a
time with a bounded step counter. Reject this branch if the path is invalid.
Count all marks and require that their total equal $c_i$.

On a successful branch, $b_u\leq R_i(u)$ for every $u$. Equality of their
finite sums forces equality at every vertex: omitting one reachable vertex
would reduce the sum by at least one. Conversely mark exactly the reachable
vertices and use, for each, the least valid bounded path code. These choices
give a successful branch. Their mathematical constructor is bounded search,
so it is primitive recursive; it need not be a logarithmic-space algorithm.
The verifier stores only a constant number of vertices and counters, never
all marks or paths.
\end{proof}

Given $c_i$, compute $c_{i+1}$ as follows. For each target $v<N$, run a fresh
exact enumeration and accumulate the bit
\[
 h_v=[v=a]\lor\bigvee_{u<N}\bigl(b_u\land E(u,v)\bigr).
\]
After the count check succeeds, this is exactly $R_{i+1}(v)$. Increment a
next-count register when it is one. The outer scan keeps only the old count,
the partial new count, and a constant number of indices. Every successful
branch produces $c_{i+1}$, and the canonical path choices show that a
successful branch exists. Iterate from $i=0$ through $N-2$.

Finally perform exact enumeration for depth $N-1$. Accept precisely when
the enumeration succeeds and the target $b$ was not marked. Cycle removal
then gives
\[
 \text{an accepting branch exists}
 \quad\Longleftrightarrow\quad
 b\text{ is not reachable from }a.
\]
This construction does not reverse nondeterministic accept and reject
states. Its complement certificate establishes an exact count and uses
that count to exclude the target.

There are $O(N^4)$ adjacency queries: at most $N$ levels, $N$ targets per
level, $N$ enumerated vertices per target, and at most $N$ edges per guessed
path. Bookkeeping has polynomial cost. Space is $O(\log N)$ plus one
adjacency check. Thus the explicit logarithmic cap construction yields
$\mathrm{NL}=\mathrm{coNL}$. Applied to polynomial-space configuration
graphs, it also gives nondeterministic polynomial-space complementation.

\paragraph{Why the induction is permitted.}
A transcript has a primitive-recursive bit-length bound because all four
loops and every path are bounded. For example, a padded format has at most
$N(N-1)+1$ enumerations; each contains $N$ mark entries, each with a length
field and at most $N$ vertex fields. With
$w=\lceil\log_2(N+2)\rceil$, its bit length is bounded by
\[
 B(N)=\bigl(N(N-1)+1\bigr)N(1+w+Nw).
\]
Use the corresponding code bound $2^{B(N)}$ for quantified transcript
variables. The stage invariant says that every surviving bounded transcript
has the correct count, and that a specified primitive-recursive canonical
transcript survives. The first assertion is bounded universal; the second
is a Boolean equality after the canonical constructor has been defined.
Both therefore compile to quantifier-free induction. The canonical
constructor uses the finite reachability table and bounded least path
search. No unbounded choice or collection is needed. A formal proof must
also show that the transcript verifier and the low-space machine describe
exactly the same bounded branches; this is a local finite realization lemma.

\subsection{True quantified Boolean formulas are PSPACE-complete}
\label{sec:tqbf}

A quantified Boolean formula has constants, Boolean connectives, and
quantifiers over the two-element domain. Its truth is a primitive-recursive
Boolean function of its finite syntax code: bounded syntax traversal and
finite assignment tables provide a direct definition. A depth-first
evaluator instead stores the current assignment and a stack of subformula
positions. For an input of length $L$ there are at most $L$ active frames,
with polynomially many bits in total. Its finite recursion tree has an
explicit exponential bound, so it can be implemented with a clock. Thus
truth of closed quantified Boolean formulas lies in PSPACE.

For hardness, consider a fixed polynomial-space machine on input $x$.
Choose fixed-width binary configuration fields. Let $d$ bits suffice for a
configuration, including a unique accepting target. Validity and adjacency
are computable by Boolean circuits of size polynomial in $|x|+s(|x|)$:
compare states and heads, check unchanged tape cells, inspect the selected
work symbols, and select the input bit by its head index. Each comparison,
selector, and local transition has an explicit finite gate construction.
There are at most $2^d$ possible codes. Invalid codes have no useful edges.

The equality-or-adjacency predicate has a polynomial-size Boolean formula
with existentially quantified auxiliary gate variables. Introduce one
variable per circuit gate, constrain it to equal that gate's Boolean
operation, and require the output bit to be one. This representation is
linear in circuit size; expanding a circuit into a formula tree would not
justify a polynomial bound. Gate values supply canonical witnesses, and
induction in gate order proves the converse.

Call this base formula $\Phi_0(a,b)$, with $a,b$ tuples of $d$ free Boolean
variables, and require endpoint validity. Recursively define
\[
\begin{split}
 \Phi_{j+1}(a,b)=\exists z\bigl(&\mathrm{Valid}(z)\land
    \forall u\,\forall v\,\bigl(
       G(a,b,z,u,v)\to\Phi_j(u,v)\bigr)\bigr),\\
 G(a,b,z,u,v)={}&(u=a\land v=z)\lor(u=z\land v=b).
\end{split}
\]
Each displayed quantifier ranges over a $d$-bit tuple, and every level uses
fresh variable names. There is only one recursive occurrence of $\Phi_j$.
Validity can be included with a polynomial-size circuit and fresh
existential gate variables at each level; its contribution remains
polynomial. If all $d$-bit codes are permitted vertices and invalid codes
are isolated, midpoint validity may instead be omitted: a path between
valid endpoints cannot benefit from an isolated invalid midpoint.

\begin{lemma}[Shared occurrence of reachability]
For valid endpoints, $\Phi_j(a,b)$ is true exactly when there is a path of
length at most $2^j$ from $a$ to $b$.
\end{lemma}
\begin{proof}
The base formula has the equality-or-adjacency meaning. If the two half-paths
exist for midpoint $z$, the induction hypothesis proves the consequent for
each of the two guarded endpoint pairs. For any other $u,v$ the Boolean
guard is false, and its implication is true. This proves the universal
formula. Conversely instantiate its universal tuples first by $(a,z)$ and
then by $(z,b)$. Both guards are true, so the two instances of the induction
hypothesis give the half-paths. Concatenate them. Every case distinction
is on a decidable finite Boolean equality.
\end{proof}

Take $j=d$. Cycle removal makes this depth sufficient, and substitute the
initial and accepting configurations. The resulting closed formula is true
exactly when the original machine accepts. The guard and quantified tuples
add $O(d)$ variable occurrences per level. The base circuit and any repeated
validity circuits have polynomial size; with $d$ levels the full formula
has polynomial size in $|x|+s(|x|)$. Binary variable identifiers add only a
logarithmic factor. A printer performs explicitly bounded nested loops over
the fields, gates, levels, and variable indices, giving a polynomial-time
many-one reduction for each fixed polynomial-space machine and its cap.

Quantifiers need not initially be in prenex form. If that convention is
required, rename bound variables, put the Boolean matrix in negation normal
form, and move each quantifier across a connective where its variable is
not free in the other operand. These finite Boolean equivalences preserve
truth and require no duplication of subformulas, giving polynomial output
size. The domain is nonempty, and all Boolean subformulas are decidable.

The iPRA invariant quantifies only over endpoint tuples below $2^d$, finite
syntax codes under the explicit size bound, and assignment codes. Truth
and repeated-squaring values are primitive-recursive Booleans. Their
equality, conjoined over bounded endpoints, is therefore an admissible
quantifier-free induction invariant after compilation. This proves the
arithmetical content of TQBF PSPACE-completeness as a family of explicit
reduction and correctness statements, without quantifying over arbitrary
sets of inputs.

\subsection{Containments and implementation obligations}

The constructions give
\[
 \mathrm{L}\subseteq\mathrm{NL}=\mathrm{coNL}
 \subseteq\mathrm{P}\subseteq\mathrm{NP}
 \subseteq\mathrm{PSPACE}=\mathrm{NPSPACE}
 \subseteq\mathrm{EXPTIME}.
\]
Deterministic computations supply single-branch nondeterministic ones.
For $\mathrm{NP}\subseteq\mathrm{PSPACE}$, enumerate all polynomial-length
witnesses with a binary counter, rerun the polynomial-time verifier, and
reuse its storage. The counter and verifier need polynomial space although
the total time is exponential. For
$\mathrm{PSPACE}\subseteq\mathrm{EXPTIME}$, the deterministic configuration
cutoff is $2^{p(n)}$ for an explicit polynomial $p$ depending on the fixed
machine. Simulating that many transitions with polynomial overhead preserves
an exponential-time bound. None of these arguments decides any of the
standard unresolved equalities between neighboring classes.

For Lean, the semantic layer may start with finite vertices and a Boolean
adjacency function. The arithmetic layer must replace this function by a
finite table or a particular coded transition computation. Required bridge
lemmas include the configuration encoding bound, transition--edge
correspondence, cycle-removal correctness, repeated-squaring correctness,
stack-evaluator equivalence and frame bound, the equal-count finite subset
lemma, counting-transcript soundness and completeness, branch realization,
and the shared-occurrence formula lemma. Finally the bounded invariants
must be translated into the chosen object-language proof calculus.

The executable reference implementation supplied with this report checks
these finite constructions on small graphs. Its transcript objects and
independent closure oracle are intentionally materialized; the Python code
is not claimed to be a logarithmic-space Turing-machine implementation.
Finite tests validate examples and detect implementation errors. They do
not replace the mathematical arguments or an eventual checked iPRA
derivation.

\section{An explicit constructive separation of polynomial and exponential time}
\label{sec:diagonal}

The foundation should support lower bounds as well as the verification of
upper bounds.  We now give a direct separation
\[
  \mathrm{P}\subsetneq\mathrm{EXPTIME}
\]
whose diagonal counterexample is a primitive-recursive function of the
proposed polynomial algorithm.  The argument deliberately uses a generous
exponential upper bound.  It does not require the sharp time hierarchy theorem,
an asymptotic domination principle, or any claim about
\(\mathrm{P}\) versus \(\mathrm{NP}\).

\subsection{The exact bounded simulation convention}

Here the source model is a deterministic one-tape machine with a fixed finite
tape alphabet containing the two input symbols and a blank.  The input word
occupies consecutive cells beginning at position zero; the head starts at
zero.  Accepting and rejecting states are absorbing.  Every binary word is
allowed as a program description: a word that does not describe a legal
transition table denotes a fixed machine that rejects immediately.  In
particular the bounded acceptance function is total for every description.
Here \(e\) ranges over binary-word codes \(e\ge1\), including malformed
machine payloads. A numerical descriptor containing 0 is normalized to
the empty description before this construction, as in the machine convention.

Write
\[
  B(e,w,t)\in\{0,1\}
\]
for acceptance of program \(e\) on input \(w\) after \(t\) iterations of its
transition function.  Absorption makes this equivalent to acceptance by time
\(t\).  It is a primitive-recursive function of all three natural-number
codes.  No assertion that the original program eventually halts is involved.
The code of a polynomial presentation is a triple \((e,c,k)\), where
\(c\geq 1\) and \(k\geq0\), with Boolean interpretation
\begin{equation}
  L_{e,c,k}(w)=B\bigl(e,w,c(|w|+1)^k\bigr).
  \label{eq:diag-presentation}
\end{equation}
Every such presentation defines a polynomial-time language, since a bounded
universal simulation incurs only a polynomial overhead.  Conversely, every
ordinary polynomial-time decision algorithm has a presentation of this form,
after enlarging its constant and degree.  A presentation that times out still
has a defined answer: its value in \eqref{eq:diag-presentation} is zero.

\begin{lemma}[A coarse physical simulation bound]
\label{lem:coarse-universal}
There is a fixed one-tape simulator and a fixed positive integer \(C\) that
compute \(B(e,w,t)\), supplied with a unary countdown \(t\), in at most
\[
  C\bigl(|e|+|w|+t+1\bigr)^4
\]
actual tape transitions.  The simulator, its correctness proof, and this
bound admit primitive-recursive descriptions and quantifier-free induction
proofs in iPRA.
\end{lemma}

\begin{proof}
First describe a simulator on a fixed finite number of virtual tapes.  One
tape holds the simulated tape; another holds the source transition table;
others hold the binary source-state code, scan markers, and the countdown.
If \(m=|e|\), a deliberately inefficient full scan of the transition table,
using repeated scans to compare state codes, costs at most a fixed multiple
of \((m+1)^2\) virtual transitions per simulated move.  The source-state code
has at most a fixed multiple of \(m+1\) bits under a conventional transition
table encoding.  A malformed table is detected by a finite parsing scan and
replaced by the rejecting program.

During \(t\) simulated moves, the simulated head can visit only the interval
\([-t,|w|+t]\).  The virtual tapes therefore occupy at most a fixed multiple
of
\(N=m+|w|+t+1\) cells, including the unary countdown.  Encode their cells on
one physical tape by a fixed number of tracks, with a marked virtual head on
each track.  A bounded number of complete sweeps reads the marked symbols,
updates them, and moves their markers.  Each virtual move costs at most a
fixed multiple of \(N\) physical moves.  Thus the main simulation costs at
most a fixed multiple of
\[
  t(m+1)^2N\leq N^4.
\]
Initialization, parsing, copying, and clearing markers fit the same bound
after increasing the fixed constant.

For formal verification, the invariant says that every encoded virtual cell
and every head marker agrees with the corresponding simulated configuration,
that the countdown equals the remaining fuel, and that the occupied interval
has the displayed bound.  All cell comparisons and scan-prefix conditions
are bounded by explicit numerical terms.  Their Boolean conjunction is a
quantifier-free induction invariant.  The running-time argument uses the
same invariant together with a transition counter.  The constant \(C\) is
fixed by the chosen finite scan implementation; this proof does not pretend
to supply an optimized numerical transition count.
\end{proof}

\subsection{A description that carries its own clock}

Let \(1^r\) denote a unary block of \(r\) ones, and let juxtaposition denote
word concatenation.  Define
\begin{equation}
  W(e,c,k)=1^{|e|}\,0\,e\,1^c\,0\,1^k\,0,
  \qquad e\ge1,\ c\geq1.
  \label{eq:self-description}
\end{equation}
The first zero terminates the unary length of \(e\).  The parser then reads
exactly that many bits as the program payload, so zeros inside \(e\) cause no
ambiguity.  The remaining two unary fields end at the next two zeros; the
word must end there.  The field for \(k=0\) consists just of its terminating
zero.  A payload need not be a legal machine description, because the default
rejecting interpretation has already made the semantics total.

The constructor \(W\) and this parser are primitive recursive.  Direct
concatenation identities prove that parsing \(W(e,c,k)\) returns exactly
\((e,c,k)\).  Its length is
\begin{equation}
  n=|W(e,c,k)|=2|e|+c+k+3.
  \label{eq:diag-length}
\end{equation}
Consequently \(|e|\leq n\), \(c\leq n\), and \(k\leq n\).

Define the Boolean language \(D\) by the following finite algorithm.  On
input \(w\), parse the displayed grammar.  If parsing fails, or if the middle
unary field has length zero, return zero.  Otherwise recover \(e,c,k\), put
\(n=|w|\), and return
\begin{equation}
  D(w)=1-B\bigl(e,w,c(n+1)^k\bigr).
  \label{eq:diag-definition}
\end{equation}
The input \(w\) itself is passed to the simulated program.  All operations
in \eqref{eq:diag-definition} are primitive recursive: parsing, length, exponentiation, bounded
iteration, and Boolean complementation.  This establishes definability in
iPRA.  A separate machine-cost proof is needed to establish the asserted
exponential upper bound.

\begin{lemma}[An exponential bound for the diagonal language]
\label{lem:diagonal-exp}
For a fixed positive integer \(K\), a fixed deterministic machine computes
\(D(w)\) in at most
\[
  2^{K(|w|+1)^2}
\]
transitions on every input \(w\).  The bound and correctness have iPRA proofs.
\end{lemma}

\begin{proof}
Consider a well-formed input of length \(n\).  By \eqref{eq:diag-length}, its recovered constants
satisfy \(1\leq c\leq n\), \(0\leq k\leq n\), and \(n\geq1\).  Set
\(T=c(n+1)^k\).  Quantifier-free induction proves \(n+1\leq2^n\) for
\(n\geq1\), whence
\begin{equation}
  T\leq n(n+1)^n
    \leq 2^n 2^{n^2}
    =2^{n^2+n}
    \leq 2^{(n+1)^2}.
  \label{eq:diag-clock-bound}
\end{equation}

An especially simple implementation produces a unary countdown of length
\(T\).  Start with the unary \(c\) block and multiply its length by \(n+1\),
\(k\) times, by repeated copying.  Every intermediate length is at most
\(T\).  In a naive tape implementation, a transferred unary mark requires
only a fixed number of scans of the input, old block, new block, and loop
markers, occupying \(O(n+T+1)\) cells.  There are at most \(n\) multiplication
stages and at most \(T\) transferred marks per stage, with polynomially
bounded loop overhead.  Thus the generous bound
\(A_0(n+T+1)^4\) covers countdown production.  The simulation lemma, with
\(|e|\leq n\), covers evaluation of \(B\) by another fixed multiple of the
same fourth power.  Parsing malformed inputs and returning zero costs only
a polynomial, also covered after increasing the constant.  For some fixed
positive integer \(A\), the complete algorithm therefore has cost
\[
  A(n+T+1)^4
\]
on well-formed inputs, with the polynomial parsing bound on the others.

Put \(q=(n+1)^2\).  Both \(T\) and \(n+1\) are at most \(2^q\), so
\[
  A(n+T+1)^4\leq 16A\,2^{4q}\leq2^{Kq}
\]
for any fixed \(K\geq4+\lceil\log_2(16A)\rceil\).  If one wishes to avoid
logarithms even in the specification of the constant, \(K=4+16A\) works.
Increase \(A\), and hence \(K\), to cover invalid inputs and the finitely many
small input lengths.  These are explicit polynomial and exponential bounds,
not an appeal to an eventual asymptotic comparison.

The iPRA audit uses Boolean prefix invariants for parsing and copying, the
simulator invariant of Lemma~\ref{lem:coarse-universal}, and the elementary
quantifier-free inductions for multiplication, exponentiation, and \eqref{eq:diag-clock-bound}.
All storage and all loop lengths are bounded by the displayed
primitive-recursive terms.  Unary countdown production is intentionally
inefficient, but its cost remains within the claimed exponential bound.
\end{proof}

\begin{theorem}[A uniform primitive-recursive mismatch witness]
\label{thm:explicit-diagonal}
iPRA proves that, for every program description \(e\), every \(c\geq1\), and
every \(k\geq0\), with description code \(e\ge1\),
\begin{equation}
  D\bigl(W(e,c,k)\bigr)
  \ne
  B\!\left(e,W(e,c,k),
       c\bigl(|W(e,c,k)|+1\bigr)^k\right).
  \label{eq:diag-mismatch}
\end{equation}
Consequently the displayed \(D\) belongs to \(\mathrm{EXPTIME}\) and has no
polynomial-time presentation.  In the intended extensional interpretation
of the coded classes, \(\mathrm{P}\subsetneq\mathrm{EXPTIME}\).
\end{theorem}

\begin{proof}
The parser equations give the exact triple \((e,c,k)\) on input
\(W(e,c,k)\).  Substituting it in \eqref{eq:diag-definition}, the left side of \eqref{eq:diag-mismatch} is the Boolean
complement of the right side.  The two possible Boolean values show that
they are unequal.  This uses no totality promise for the original program:
both rejecting and nonhalting computations have bounded-acceptance value
zero, whose complement is one.

Lemma~\ref{lem:diagonal-exp} supplies the exponential-time upper bound.
Equation \eqref{eq:diag-mismatch} rules out each polynomial presentation, with the actual
counterexample word given by the same primitive-recursive constructor
\(W\).  If an ordinary polynomial-time decider computed \(D\), choose a
clock \(c(n+1)^k\) dominating its running time.  Its bounded-acceptance
value would equal \(D\) on every word, contradicting \eqref{eq:diag-mismatch} on
\(W(e,c,k)\). For the non-strict inclusion, replace each polynomial
descriptor by the total clocked evaluator that returns its Boolean answer,
including its timeout answer. That evaluator has a genuine polynomial
runtime bound by the simulator lemma. An exponential clock dominates this
bound and therefore preserves its answer. Merely increasing the clock of
the original possibly partial program would not justify the inclusion.
\end{proof}

\subsection{What the separation does, and what its formalization must retain}

Equation \eqref{eq:diag-mismatch} is stronger, constructively, than merely negating the existence
of a polynomial decider.  Intuitionistic logic does not in general turn
\(\neg\forall w\,R(w)\) into \(\exists w\,\neg R(w)\).  Here there is no such
conversion: \(W(e,c,k)\) explicitly gives a differing input for each proposed
presentation.  The statement can be expressed as the universal closure of
a quantifier-free Boolean inequality with named primitive-recursive
functions.  The required constructor/parser identities are themselves
proved by quantifier-free induction on finite word operations.

The unary clock fields are essential to the resource estimate.  If the
exponent \(k\) were given in unrestricted binary and then used without a
length-dependent cap, it could be exponential in \(|w|\), and the same
algorithm could require more than exponential time.  Unary encoding still
represents every fixed polynomial clock, while guaranteeing \(k,c\leq|w|\)
on the specially constructed inputs and every successfully parsed input.

The argument also keeps program evaluation and logical reflection separate.
It evaluates a specified machine for a specified finite number of steps;
it does not evaluate arbitrary primitive-recursive definitions uniformly,
decide whether an arbitrary theorem is true, or infer soundness from an
arbitrary iPRA proof code.  The familiar diagonal phenomenon therefore poses
no conflict with the totality of the bounded evaluator.

For a Lean formalization, the key interfaces are the word constructor,
its exact parsing theorem, the bounded evaluator, its Boolean-valuedness,
the simulator's physical transition bound, and the elementary exponential
inequalities.  The final mismatch proof is short once those interfaces have
been established.  A checked Lean proof of the mismatch algebra alone does
not establish the machine-time bound or automatically produce an iPRA
derivation.  The completed development must retain the finite invariants and
the code-generating bridge described above for those separate obligations.


\section{Circuits, counting, alternation, and finite probability}
\label{sec:extensions}

These extensions illustrate that the foundation is reusable. They involve
finite sums, products and syntax, with resource bounds proved separately.
They require neither infinite probability spaces nor quantification over an
actual infinite family as a new sort.

\subsection{Uniform and nonuniform circuits}

For a Boolean presentation \(A(x)\), a polynomial-size circuit presentation
means that for each length \(n\) there exists a circuit \(C_n\) of at most
\(p(n)\) gates agreeing with \(A\) on all \(2^n\) words of that length.
This is the ordinary nonuniform class P/poly, read schematically over \(A\)
or arithmetically when \(A\) is a supplied Boolean numerical function.
There is no automatic efficient operation that selects \(C_n\) from an
existence statement.

If a fixed polynomial-time machine computes \(A\), unroll its computation
with all \(n\) input bits left as circuit inputs, instead of hardwiring
\(x\) as in Cook--Levin. The same cell induction proves correctness, and
the number of gates is polynomial in \(n\). Moreover a generator constructs
\(C_n\) in polynomial time from \(1^n\). Conversely, if a polynomial-time
generator on \(1^n\) produces correct polynomial-size circuits, generating
and evaluating the circuit is a polynomial-time decision procedure. Thus
\ipra{} supports both the nonuniform upper-bound construction and the
equivalence with this particular polynomial-time uniformity convention.
The unary presentation of \(n\) is part of that convention.

\begin{theorem}[Constructively selected hard truth tables]
\label{thm:circuitcount}
There is a primitive-recursive function \(H\) such that, for every \(n\ge4\),
\(H(n)<2^{2^n}\) is the truth table of a Boolean function not computed by
any \(n\)-input circuit of at most
\[
 s(n)=\left\lfloor\frac{2^n}{16n}\right\rfloor
\]
internal AND, OR, or NOT gates. Constants and inputs are available as wires.
For each such circuit a primitive-recursive bounded search returns an input
on which it disagrees with this truth table. These are \ipra{} theorems.
\end{theorem}

\begin{proof}
Put \(s=s(n)\) and \(R=n+s+2\), counting two constant wires. Pad a circuit
with unused gates to exactly \(s\) gates; this preserves its function.
An overcount of the exactly-\(s\) descriptions is
\[
 K=R(3R^2)^s:
\]
choose an output wire, and for each gate one of three operations and two
predecessor indices. Ignore invalid forward references; doing so only
increases the count. NOT gates simply ignore their second index. Thus
\(K\) bounds the number of \emph{functions} represented using at most
\(s\) gates. We do not claim that padding is injective on all shorter
syntactic descriptions.

For \(n\ge4\), \(n+2\le2^{n-1}\) and \(s\le2^{n-1}\), hence
\(R\le2^n\). Also \(3R^2\le2^{4n}\), so
\[
 K\le2^{n+4ns},\qquad
 n+4ns\le n+2^{n-2}\le2^{n-1}<2^n.
\]
The inequality \(n\le2^{n-2}\) starts at \(n=4\) and is preserved by
successor; the other inequalities follow by multiplication and integer
division. Hence \(K<2^{2^n}\), the number of possible truth tables.

Enumerate the \(K\) candidate descriptions, mapping invalid ones to a
fixed table, and compute each valid circuit's \(2^n\) output bits. Maintain
an image bitset indexed by the \(2^{2^n}\) possible truth tables. After
processing \(i\) descriptions, its number of set bits is at most \(i\):
inserting one value adds at most one new mark. This invariant uses bounded
sums and Boolean comparisons, so quantifier-free induction proves it.
Since the image cannot cover the table universe, bounded least search
returns a missing table. Define this table to be \(H(n)\), taking a fixed
default value for \(n<4\). Every loop and the image code itself have an
explicit primitive-recursive bound, even though they are enormous.

For a candidate circuit, equality with \(H(n)\) at all inputs would make
its truth-table integer equal to \(H(n)\), contrary to missingness.
More constructively, scan the \(2^n\) positions for the first differing
bit. If the scan had no success, its finite all-equal bit would be 1 and
the two table integers would be equal. Thus the scan supplies a mismatch
without turning an unbounded negation into an existence claim.
\end{proof}

This produces a primitive-recursive hard family, and therefore a
primitive-recursive language obtained by consulting the appropriate table
bit. Its construction can take far more than exponential time. It is not
an exponential circuit lower bound for NP, P, or another small uniform
complexity class. The example separates constructive existence from
efficient uniform construction with no ambiguity.

\subsection{Counting witnesses}

For a fixed polynomially bounded verifier, define
\[
 \#V(x)=\sum_{j<2^{p(n)}}
       \Acc_{\#}(e,x,\operatorname{bits}_{p(n)}(j),q(n)).
\]
This is a primitive-recursive natural-valued function and
\(0\le\#V(x)\le2^{p(n)}\). Its output has at most \(p(n)+1\) bits.
These presentations give the function class \(\#\mathsf P\).
Fixing exact certificate lengths makes the counted objects unambiguous;
padding conventions must be specified because they can change the count.

The support identity
\(\#V(x)>0\leftrightarrow N_v(x)=1\) follows from bounded witness search.
Disjoint tagged certificates add counts; pairs of certificates multiply
counts, by finite distributivity. Thus addition and multiplication of two
such functions have explicit verifier constructions after padding without
duplicating accepted certificates. For addition, use a tag and require all
padding bits to be zero; for multiplication, concatenate the two certificate
fields. This extra padding check is essential for exact counting.

When the circuit extension in Cook--Levin is unique for each input
assignment, counting satisfying gate assignments with the asserted output
counts exactly the accepting witnesses. This is a parsimonious statement
for that specified fixed-width witness representation. It does not assume
that an arbitrary many-one reduction preserves numbers of solutions.

\subsection{A fixed polynomial hierarchy}

For each fixed \(i\), consider presentations
\[
 \exists y_1<2^{p_1(n)}\ \forall y_2<2^{p_2(n)}\ \cdots\ Q_i y_i<2^{p_i(n)}
 \quad V(x,y_1,\ldots,y_i)=1,
\]
where the number of blocks \(i\) is fixed, each \(p_j\) is an explicit
polynomial, and \(V\) is a polynomially clocked Boolean verifier. These
are the bounded-quantifier presentations of \(\Sigma_i^P\); starting with
a universal block gives \(\Pi_i^P\). Every individual input formula is
bounded and therefore has a primitive-recursive characteristic bit.
Complement swaps the finite quantifiers and negates the matrix; this is
valid constructively because all the underlying finite tests are decidable.

A depth-first evaluation keeps one polynomial-length assignment for each
of the fixed number of blocks, the verifier workspace, and loop counters.
It therefore uses polynomial space. Its total number of branches is at
most \(2^{p_1(n)+\cdots+p_i(n)}\), with a polynomial factor for verification.
The resulting compiler proves that every fixed level lies in PSPACE.
Varying \(i\) as a metatheoretic family describes PH. This does not assert
that all levels collapse, nor that a single polynomial bound works for
arbitrary variable-level formulas. Section~\ref{sec:tqbf} handles a
variable-length quantified formula as its own finite input.

\subsection{Randomness is finite integer counting}

A randomized descriptor supplies a deterministic clocked machine, a
polynomial number \(p(n)\) of random bits, and a polynomial runtime.
Random bits are simply an additional exactly sized input field; early
termination ignores unused bits. For each input let
\[
 M=2^{p(n)},\qquad
 C(x)=\sum_{j<M}\operatorname{out}(x,\operatorname{bits}_{p(n)}(j)).
\]
The acceptance probability notation \(C(x)/M\) abbreviates a finite
rational number. Inequalities are represented by cross multiplication of
integers, so no measure theory or real numbers are needed.

There is a constructive subtlety when the intended language is an arbitrary
formula \(A(x)\). The two implications
\[
 A(x)\to 3C(x)\ge2M,\qquad
 \neg A(x)\to 3C(x)\le M
\]
alone need not identify \(A\) constructively: a high count gives only
\(\neg\neg A\). We therefore specify a BPP presentation by a proved gap
\[
 \forall x\ge1\quad
       (3C(x)\le M)\ \lor\ (3C(x)\ge2M)
\]
and pointwise equivalence of \(A(x)\) with the finite majority bit
\([2C(x)\ge M]\). Equivalently one can supply a Boolean characteristic
\(\chi_A\) and prove that the number of wrong-output seeds is at most
\(M/3\). The characteristic is not assumed to be computable in polynomial
time. In the standard extensional reading this is the usual BPP notion,
with its constructive identification made explicit.

\begin{theorem}[Integer error amplification]
\label{thm:amplification}
Suppose a randomized Boolean algorithm has at most \(F\) failing seeds out
of \(M\), where \(3F\le M\). The majority of
\(r=16(k+1)+1\) independent runs has at most a \(2^{-k}\) fraction of
failing combined seeds. The construction and the integer inequality are
provable in \ipra. For a polynomial-time randomized descriptor its runtime
and random-bit usage are polynomial in the input length and \(k\).
\end{theorem}

\begin{proof}
Write \(r=2m+1\) and consider all \(M^r\) tuples of seeds. Let \(S(\rho)\)
be the number of failing runs on tuple \(\rho\), and let \(F_0\le F\)
be the actual number of failing single-run seeds. Each successful seed
contributes weight 1 and each failing seed weight 2. Finite product
expansion gives
\[
 \sum_{\rho}2^{S(\rho)}=(M+F_0)^r\le(M+F)^r.
\]
If \(B\) tuples have a failing majority, each has weight at least
\(2^{m+1}\). Since \(3(M+F)\le4M\),
\[
 2^{m+1}B\le(M+F)^{2m+1},\qquad
 3\cdot9^m B\le2\cdot8^m M^{2m+1}.
\]
The second inequality follows by multiplying the first by
\(3^{2m+1}\), using the preceding bound, and canceling a positive common
power of two. All these are integer operations.

The concrete inequality \(2\cdot8^8=33554432<43046721=9^8\) gives, for
\(m=8(k+1)\),
\(2^{k+1}8^m<9^m\). Combining the inequalities yields
\(2^kB\le M^r\), the asserted bound. The proof of the product expansion
is induction over the number of independent seed fields; its invariant is
an equality of finite sums. Inequality preservation and counting bad
tuples are equally bounded. Thus only quantifier-free induction is used
after compilation.

Implement the independent runs sequentially, retaining a count of output
bits. The combined random input has \(r p(n)\) bits, and the runtime is
a fixed polynomial majorant of \(r\) times the one-run time plus input
and counter bookkeeping. This is polynomial in \(n,k\). No extraction of
a single good seed for all inputs, and no derandomization theorem, is
claimed.
\end{proof}

Deterministic machines ignore their random bits, giving
\(\Pclass\subseteq\mathsf{BPP}\). Exhaustive seed counting uses polynomial
space: retain one \(p(n)\)-bit seed, a \(p(n)+1\)-bit acceptance counter,
and the polynomial verifier workspace. Therefore
\(\mathsf{BPP}\subseteq\PSPACE\). A one-sided RP presentation has no
accepting seeds on no-instances and at least half its seeds accepting on
yes-instances, together with the same explicit language identification.
A successful seed is an NP witness, giving \(\mathsf{RP}\subseteq\NP\).
Complement gives \(\mathsf{coRP}\subseteq\coNP\). These arguments do not
assume that BPP equals P or that randomized expected-time bounds can always
be replaced by deterministic clocks with identical behavior.

\subsection{Worst case, finite averages, and amortized costs}

For a supplied total clocked algorithm, its actual stopping time below the
clock is a bounded least search. Maximum and minimum over the \(2^n\)
inputs are primitive recursive. For an explicitly given finite rational
distribution, multiply costs by its integer weights and divide by their
positive total weight. The average lies between the minimum and maximum
by summing the corresponding pointwise inequalities.

Amortized analysis concerns a different finite quantity. If states
\(s_0,\ldots,s_m\), nonnegative costs \(c_i\), and a numerical potential
\(\Phi\) obey
\(c_i+\Phi(s_{i+1})\le a_i+\Phi(s_i)\), summing and canceling gives
\[
 \sum_{i<m}c_i+\Phi(s_m)
       \le\sum_{i<m}a_i+\Phi(s_0).
\]
The induction is an inequality between finite sums. This provides the
usual amortized upper bound when \(\Phi\ge0\), and does not posit an
ordering of amortized and average-case complexity without additional
assumptions about workloads and distributions.

\section{Lean formalization: checked core and complete roadmap}
\label{sec:lean}

\subsection{What has actually been checked}

The standalone project in \code{lean/} was compiled with Lean 4.19.0 and
uses only the bundled \code{Std} library. It contains 31 named theorems,
six audited proof-bearing definitions, and 12 checked examples. There are
no \code{sorry} placeholders or added axiom declarations. The recorded
dependency audit contains no \code{sorryAx} or \code{Classical.choice};
some proofs depend on Lean's standard \code{propext} and \code{Quot.sound}.
These dependencies and the complete build output are supplied in
\code{verification/lean-build-and-axioms.log}.

Lean's axiom audit reports dependencies of the elaborated statement and
proof \cite{LeanAxioms}. Even an axiom-free ambient theorem would not by
itself prove that its statement has a derivation in first-order \ipra.
Unrestricted Lean induction, arbitrary function parameters and richer types
are part of the metatheory until translated. We explicitly retain that
distinction.

\begin{center}
\begin{tabularx}{\textwidth}{@{}p{0.23\textwidth}X@{}}
\toprule
Module & Checked content and its boundary\\
\midrule
\code{Bounded} & Boolean finite existential/universal semantics, failed
search semantics, witness soundness and completeness. The mathematical
leastness argument is not a separate checked theorem.\\
\code{Machine} & Iteration composition, fixed points, accepting-by-time
monotonicity and trace uniqueness for an abstract fixed transition system.
It is not yet an encoded Turing-machine interpreter.\\
\code{Reductions} & Nonnegative polynomial syntax, evaluation, substitution,
monotonicity, reduction composition, and size/call-cost certificate
composition. Supplied costs still need a machine implementation proof.\\
\code{Reachability} & An independent directed-walk relation, splitting and
concatenation, stutter removal, and the exact midpoint-recurrence theorem
for paths of length at most \(2^k\). No Savitch storage bound is inferred
from this semantic result alone.\\
\code{Circuits} & AND/NOT/constant Tseitin constraints, correctness for a
whole gate system, width and linear clause-count bounds. The general
machine-to-circuit compiler is not yet a checked Lean theorem.\\
\bottomrule
\end{tabularx}
\end{center}

The central checked reachability statement, in the supplied namespace, is
\code{reach\_iff\_short\_walk}. It relates the executable Boolean recurrence
to a separately defined inductive walk relation, including bounded
endpoints. This avoids validating the recurrence merely by restating its
definition. The Tseitin constraints similarly compare finite clause
evaluation with gate equations, not with an assumed SAT specification.

\subsection{Definitions needed for the complete numerical layer}

The report fixes the mathematical meaning of each definition. The remaining
coding work should produce primitive-recursive syntax, not simply more
ambient functions. The following interface list is sufficient to organize
the implementation; all arguments and results are natural-number codes,
and predicates return bits.

\begin{longtable}{@{}>{\RaggedRight\arraybackslash}p{0.29\textwidth}>{\RaggedRight\arraybackslash}p{0.63\textwidth}@{}}
\toprule
Numerical operation & Required definition and invariants\\
\midrule
\endhead
Words and records & Validity, length, bit access, concatenation, slicing,
mixed-radix arrays, update, and decode/encode equations. Word code 0 is
reserved; language input theorems are guarded by \(x\ge1\).\\
Machine coding & Fixed global alphabet; finite transition-table parsing;
state labels; malformed-code defaults; signed head and finite tape codes;
\(\operatorname{Init},\Step,\Run,\Acc,\Acc_{\#}\).\\
Trace checker & Window conversion; explicit configuration bound;
initial-row test; bounded conjunction of adjacent-step equations; final
state and resource checks.\\
Descriptors & Polynomial and exponential clock evaluation; normalization
of invalid descriptors; witness length; \(P_d,N_v\); output extraction and
clocked transducers.\\
Circuit syntax & Gate-list encoding, predecessor validity, input and output
counts, prefix evaluation, CNF list encoding, literal evaluation, and
assignment validity.\\
Compilers & Code-level machine composition; local unrolling;
\(F_v(x),\operatorname{Ext}_v,\operatorname{Restr}_v\); source and output
bit-length bounds and target clocks.\\
Configuration graph & \(N_0(e,n,s)\), mixed-radix projections, cap and
transition tests, initial and target codes, adjacency.\\
Graph algorithms & Bounded path coding, split/join, cycle removal, linear
marking, repeated-squaring table updates, clocked stack evaluator,
exact-count transcript checking and canonical transcript generation.\\
QBF syntax & Finite Boolean terms and quantifiers, capture-free variable
renaming, guarded shared-occurrence constructor, bounded truth evaluation,
and its polynomial-space evaluator.\\
Diagonal syntax & \(W(e,c,k)\), exact parser, the clocked complemented
simulation, unary clock construction, mismatch equation.\\
Finite counting & Sum/product folds, weighted counts, image bitsets,
bounded least missing element, seed-tuples and majority.\\
\bottomrule
\end{longtable}

For an arbitrary graph use an adjacency bit table of length \(N^2\).
For an implicit configuration graph use the particular primitive-recursive
transition test with numerical machine parameters. The latter does not
require evaluation of an arbitrary PR program supplied as input. Similarly,
QBF truth has a fixed finite grammar with bounded Boolean values; it is not
truth for arbitrary first-order arithmetic syntax.

\subsection{The actual object-theory target in ProveIt}

At the pinned revision, \code{LeanProofs.ConciseIPRA} supplies
\code{PR}, \code{PTm}, \code{Fm}, \code{PRAAx}, and the intuitionistic
proof judgment \code{Prf}. The PR constructors are zero, successor,
projection, composition and primitive recursion. The desired endpoint for
a closed sentence \(\varphi\) is the Lean type
\[
  []\ \vdash[\mathsf{PRAAx}]\ \varphi,
\]
written in the repository as
\code{([] : List (Fm PTm 0))} followed by its
\code{\(\vdash\)[PRAAx]} notation and the encoded formula.
This is a type of actual derivations, not a theorem saying that the
corresponding proposition happens to hold in Lean's natural numbers.

The inspected interfaces include:
\begin{itemize}
\item \code{PRA.comp}, \code{PRA.precZero}, \code{PRA.precSucc} for
primitive-recursive defining equations, and \code{PRA.qfInd} for the
permitted induction rule.
\item \code{BoundedFormulas.bf\_reflect}, \code{bf\_dec},
\code{bf\_choice}, and \code{bf\_least}, for fixed bounded syntax and its
numeric Boolean interpretation.
\item \code{LeastBoundedWitness.found\_iff}, \code{sentinel\_iff},
\code{decision}, and \code{skolem}, for bounded numerical witnesses.
\item The \code{BoundedSyntaxInduction} namespace provides
\code{induction}, \code{induction\_at},
\code{induction\_rule}, and \code{bounded\_induction\_closed}. Their
source applies \code{PRA.qfInd} to the compiled-bit equation and translates
back by reflection of that fixed bounded syntax.
\end{itemize}

The exact namespaces and source paths are recorded in
\code{lean/OBJECT\_IPRA\_ROADMAP.md} and \code{SOURCE\_AUDIT.md}.
The pinned repository uses Lean 4.32.0 and its own Mathlib revision. This
package's independent Lean 4.19.0 core is not claimed to have been tested
as a drop-in module under that different toolchain. Integration should use
the repository's pinned environment and port or restate the small semantic
lemmas there; it should not overwrite the upstream toolchain.

\subsection{A staged completion plan with concrete proof obligations}

\paragraph{Stage 1: code the arithmetic primitives.}
For each function in the interface table construct a term of the actual
\code{PR k} syntax. Start with word and array operations, bounded folds,
polynomials and finite transition tables. Derive their equations using the
existing ordinary arithmetic library. Prove encode/decode and bound lemmas
before invoking the semantic companion. In particular, conversion from
sparse tapes to fixed windows is a proved map between two encodings.

\paragraph{Stage 2: transport the bounded kernel.}
Translate \code{bany\_correct}, \code{ball\_correct}, and witness
soundness/completeness to numeric 0/1 bits. Reuse the existing bounded-syntax
reflection endpoints instead of defining a second logical kernel. Translate
iteration and trace uniqueness using the invariant
\[
 \forall i\le\min(j,T)\quad
 \operatorname{row}(z,i)=\Run(e,x,i),
\]
with the explicit row and trace bounds. Apply the existing bounded-induction
endpoint to its fixed bounded syntax. No arbitrary function-valued trace
survives in the object language.

\paragraph{Stage 3: implement and certify costs.}
Choose the finite machine implementations for universal simulation, input
copying, polynomial clock construction, sequential composition, and output
printing. Give a transition counter and a storage-height counter. Prove
that their numeric operational traces meet the cost polynomials. This
connects a \code{PolyCost} certificate to actual machine transitions;
the current algebraic certificate alone does not make that connection.

\paragraph{Stage 4: complete Cook--Levin.}
Prove the local update truth table implements one machine move. Then prove
the row invariant, topological validity of emitted wires, output size and
generator runtime. Add acyclic evaluation and unique extension to the
already checked Tseitin-gadget core. The final endpoints are the two
witness implications of Theorem~\ref{thm:cook}, followed by existential
introduction/elimination and explicit target-clock normalization. A useful
integration target is an object-theory theorem with \(v,x,y\) free, later
universally closed, rather than a type quantifying over arbitrary Lean
languages.

\paragraph{Stage 5: complete finite-graph and space proofs.}
Replace ambient inductive walks by bounded path checkers. For repeated
squaring through a fixed outer level \(K\), bound every path code by
\(\max(2,N)^{2^K+1}\) and every table by \(2^{N^2}\). Translate the
checked walk splitting and midpoint theorem. Separately verify the
iterative stack machine, with at most \(K+1\) frames and its explicit
time clock. For complementing nondeterminism, implement both the transcript
verifier and the streaming branch machine, and prove branch realization in
both directions. The exact-count invariant uses the transcript bound in
Section~\ref{sec:space}. For QBF, prove capture-free renaming and the
shared-occurrence formula lemma, then the printer's size/time bounds.

\paragraph{Stage 6: lower the diagonal and extensions.}
The parser equations and Boolean mismatch need only a small object proof
once word operations are available. The resource membership additionally
needs Stage 3's real simulator and unary-clock bounds. Finite circuit
counting uses bounded image cardinality; amplification uses finite product
expansion. Each receives its own bounded induction invariant and output
sentence. These are separate milestones, so the difficult implementation
lemmas are neither assumed as axioms nor hidden behind a generic
``formalizable'' label.

At each stage the final object proof can be checked by the ordinary
\code{Prf PRAAx} interface. If a native scalar-core export is desired,
apply the repository's proved ordinary/native equivalence only after the
ordinary derivation exists. There is no need to introduce a new foundation
for each complexity result.

\section{Induction and verification ledger}
\label{sec:ledger}

\subsection{The induction audit}

The following table records the actual finite information used by the main
proofs. Bounds may be intentionally much larger than the running program's
storage, because proof coding and resource implementation are distinct.

\begin{longtable}{@{}>{\RaggedRight\arraybackslash}p{0.23\textwidth}>{\RaggedRight\arraybackslash}p{0.35\textwidth}>{\RaggedRight\arraybackslash}p{0.34\textwidth}@{}}
\toprule
Result & Induction invariant & Numerical bounds / witnesses\\
\midrule
\endhead
Bounded search & Found-bit and saved witness agree with the scanned prefix;
all earlier candidates fail. & Scan length \(r\); index at most \(r\);
Boolean prefix products.\\
Bounded simulation & Tape cells, head and state match the machine; time
counter advances once per move. & Window \(n+2T+3\), row bound
\(qWg^W\), explicit trace bound.\\
Cook--Levin & Every label in the processed row equals the corresponding
run label; all wire predecessors are earlier. & \(t\le T\), \(j<W\),
gate count \(a_v(T+1)W\); evaluation is the extension witness.\\
Tseitin converse & All processed gate variables equal the evaluation
from the input prefix. & Finite gate/assignment array; prefix projection
is the witness map.\\
Cycle removal & Endpoints and adjacency preserved; repeated-vertex
deletion decreases length. & At most the original path length many
deletions; least repeated pair.\\
Savitch & The finite table or stack answer equals a path of length
at most \(2^j\). & Path bound \(N^{2^K+1}\) at outer level \(K\);
stack depth \(K+1\); explicit recursion-node clock.\\
Space complement & All surviving transcripts have the true count; a
specified canonical transcript survives. & Transcript code below
\(2^{B(N)}\); paths chosen by bounded least search.\\
QBF completeness & Guarded formula truth equals the bounded reachability
bit for every endpoint pair. & Endpoints below \(2^d\); formula and
assignment bounds; fresh variable indices.\\
Diagonal separation & Parser returns the supplied fields; arithmetic
clock inequalities; simulator-prefix agreement. & Unary \(c,k\le n\);
witness \(W(e,c,k)\); clock \(2^{K(n+1)^2}\).\\
Circuit lower bound & The image of \(i\) circuit descriptions has at most
\(i\) marked truth tables. & Finite table universe \(2^{2^n}\); least
missing table and least disagreeing input.\\
Amplification & Weighted tuple count factors as a product; majority-bad
tuples have minimum weight. & \(M^r\) seed tuples; integer powers;
no infinite probability object.\\
\bottomrule
\end{longtable}

\subsection{Executable checks supplied with the archive}

The graph implementation was compared against an independent
Floyd--Warshall closure on all 530 directed graphs, including self-loops,
with one, two or three vertices, covering 4,674 source--target pairs.
Linear reachability, midpoint reachability, bounded path search, and
canonical nonreachability certificates agreed. Separate tests check the
equal-count argument and reject damaged paths, marks, counts and transcripts.
The reference implementation materializes tables and certificates; it is
not claimed to run in logarithmic space.

The compiler implementation performs an actual finite single-tape
simulation-to-circuit-to-3CNF construction. Its reference convention uses
a fixed binary prefix followed by a fixed-width certificate; the general
report permits a separator as well. Tests compare the independent simulator
and circuit at every relevant row and tape cell, include leftward moves,
writes, empty inputs, clock zero and terminal absorption, and validate
explicit assignment extension and restriction. They include all 86 full
assignments of seven tiny circuit examples, another 32 assignments covering
unused inputs and both constant outputs, 32 aliased-AND assignments,
42 parity-compiler instances, and 107 witness/clock cases. Independently
flipping each of 2,419 auxiliary variables in one compiled example is
rejected by its gate constraints.

The graph suite has five tests and the circuit suite eleven, all passing in
the captured run. These are finite validation evidence for the executable
artifacts. They are not a substitute for the general mathematical proofs,
nor a claim of completed object-theory formalization.

\subsection{Boundary statements and next research directions}

The foundation does not settle P versus NP, NP versus coNP, L versus NL,
or PSPACE versus EXPTIME. It also does not turn arbitrary extensional
correctness into a short proof, infer machine costs from PR definability,
or reflect arbitrary \ipra{} proof codes into truth. Sharp hierarchy
theorems need explicit constructibility and simulation hypotheses beyond
the coarse separation proved here. Uniform randomized or quantum
complexity needs its own representation and error conventions. Completeness
or speedup results that quantify over arbitrary total algorithms need a
separate analysis of which clocks and totality proofs can be represented.

The next concrete integration milestone is to lower the bounded-machine
and circuit-construction lemmas to the existing derivation library. That
would deliver an actual object-theory Cook--Levin theorem without first
formalizing every part of the subject. The graph kernel and the inspected
bounded-induction endpoints then provide a second path to Savitch and
space complementation. This modular order makes every claimed gain
checkable and avoids adding unnecessary axioms to \ipra.

\begin{thebibliography}{20}
\bibitem{WikiComplexity}
Wikipedia contributors. \emph{Computational complexity theory}.
Subject map requested by the user; consulted 20 November 2026.
\url{https://en.wikipedia.org/wiki/Computational_complexity_theory}.

\bibitem{Cook1971}
Stephen A. Cook. The complexity of theorem-proving procedures.
\emph{Proceedings of the Third Annual ACM Symposium on Theory of Computing}
(1971), 151--158. DOI: \href{https://doi.org/10.1145/800157.805047}{10.1145/800157.805047}.
Author-hosted copy:
\url{https://www.cs.toronto.edu/~sacook/homepage/1971.pdf}.
The iPRA witness and induction audits in this report are our reconstruction.

\bibitem{CookPNP}
Stephen A. Cook. \emph{The P versus NP problem}. Millennium problem
description, 2000.
\url{https://www.cs.toronto.edu/~toni/Courses/Complexity/handouts/cook-clay.pdf}.

\bibitem{Savitch1970}
Walter J. Savitch. Relationships between nondeterministic and deterministic
tape complexities. \emph{Journal of Computer and System Sciences}
4(2) (1970), 177--192.
\href{https://doi.org/10.1016/S0022-0000(70)80006-X}{doi:10.1016/S0022-0000(70)80006-X}.
The publisher abstract was retrieved; the full publisher page denied direct
access. The proof and arithmetic audit here are self-contained.

\bibitem{Immerman1988}
Neil Immerman. Nondeterministic space is closed under complementation.
\emph{SIAM Journal on Computing} 17(5) (1988), 935--938.
Author-hosted copy: \url{https://people.cs.umass.edu/~immerman/pub/space.pdf}.

\bibitem{AroraBarak}
Sanjeev Arora and Boaz Barak. \emph{Computational Complexity: A Modern
Approach}. Cambridge University Press, 2009.
Author's background and draft page:
\url{https://theory.cs.princeton.edu/complexity/}.

\bibitem{Buss1986}
Samuel R. Buss. The polynomial hierarchy and intuitionistic bounded
arithmetic. \emph{Structure in Complexity}, Lecture Notes in Computer
Science 223 (1986), 77--103.
Author's description and article links:
\url{https://mathweb.ucsd.edu/~sbuss/ResearchWeb/intuitionisticBA/}.
This source concerns intuitionistic bounded arithmetic, not an automatic
identification of iPRA-definable functions with polynomial time.

\bibitem{ClayPNP}
Clay Mathematics Institute. \emph{P vs NP}, Millennium Problems.
Listed as unsolved when consulted 20 November 2026.
\url{https://www.claymath.org/millennium/p-vs-np/}.

\bibitem{LeanAxioms}
Lean language reference. \emph{Axioms}; \emph{Validating a Lean Proof}.
\url{https://lean-lang.org/doc/reference/latest/Axioms/}.
The companion project is pinned to Lean 4.19.0; the reference explains
dependency auditing rather than the object-iPRA translation.

\bibitem{ProveItPR}
ProveIt contributors. \emph{ConciseIPRA/Lean/ConciseIPRA/PR.lean},
with \emph{Syntax.lean} and \emph{Deduction.lean}.
Repository commit \texttt{83421593771f877bc1ae98f4229eb35507bff05b},
inspected 20 November 2026.
\url{https://github.com/VladimirReshetnikov/ProveIt/blob/83421593771f877bc1ae98f4229eb35507bff05b/Logic/PrimitiveRecursiveArithmetic/ConciseIPRA/Lean/ConciseIPRA/PR.lean}.

\bibitem{ProveItBounded}
ProveIt contributors. \emph{IPRADerivations/Foundations/BoundedFormulas.lean},
same commit. Object-theory bounded Boolean reflection and witness endpoints.
\url{https://github.com/VladimirReshetnikov/ProveIt/blob/83421593771f877bc1ae98f4229eb35507bff05b/Logic/PrimitiveRecursiveArithmetic/Derivations/Lean/IPRADerivations/Foundations/BoundedFormulas.lean}.

\bibitem{ProveItLeast}
ProveIt contributors. \emph{IPRADerivations/Foundations/LeastBoundedWitness.lean},
same commit.
\url{https://github.com/VladimirReshetnikov/ProveIt/blob/83421593771f877bc1ae98f4229eb35507bff05b/Logic/PrimitiveRecursiveArithmetic/Derivations/Lean/IPRADerivations/Foundations/LeastBoundedWitness.lean}.

\bibitem{ProveItInd}
ProveIt contributors. \emph{IPRADerivations/Foundations/BoundedSyntaxInduction.lean},
same commit. Fixed bounded syntax is lowered to quantifier-free induction.
\url{https://github.com/VladimirReshetnikov/ProveIt/blob/83421593771f877bc1ae98f4229eb35507bff05b/Logic/PrimitiveRecursiveArithmetic/Derivations/Lean/IPRADerivations/Foundations/BoundedSyntaxInduction.lean}.
\end{thebibliography}

\end{document}
